% Recovered from the previously completed PDF; no new mathematical content.
% Prose is editable; formula boxes retain original symbol positions.
% Build from repository root: xelatex -no-shell-escape -halt-on-error -output-directory=... print/fudan/NAME.tex
\documentclass[11pt,a4paper]{article}
\usepackage[margin=23mm,headheight=14pt,headsep=8.5mm,footskip=10mm]{geometry}
\usepackage{amsmath,amssymb}
\DeclareMathAlphabet{\mathsf}{OT1}{cmss}{m}{n}
\usepackage[no-math]{fontspec}
\usepackage{xeCJK}
\setmainfont{Latin Modern Roman}
\setCJKmainfont{FandolSong-Regular.otf}[Path=assets/fonts/fandol/,BoldFont=FandolSong-Bold.otf,ItalicFont=FandolSong-Regular.otf,Script=CustomDefault]
\xeCJKsetup{PunctStyle=quanjiao,CheckSingle=true}
\usepackage{setspace}\setstretch{1.12}
\setlength{\parindent}{2em}\setlength{\parskip}{3pt}
\emergencystretch=6em
\hfuzz=9pt
\usepackage{titlesec}
\titleformat{\section}{\centering\normalfont\Large\bfseries}{}{0pt}{}
\titlespacing*{\section}{0pt}{1.8\baselineskip}{1.0\baselineskip}
\usepackage{fancyhdr}\pagestyle{fancy}\fancyhf{}
\fancyhead[L]{\small 数学手记·高等代数}\fancyhead[R]{\small 22S09–22S15}
\fancyfoot[C]{\thepage}\renewcommand{\headrulewidth}{0.4pt}
\usepackage[hidelinks,unicode]{hyperref}
\hypersetup{pdftitle={22S09–22S15：题目、解答与小结},pdfauthor={Vincent Arqain}}
\newcommand{\RB}[4]{\raisebox{0pt}[#2bp][#3bp]{\hbox to #1bp{#4\hss}}}
\newcommand{\RG}[5]{\hbox to 0pt{\kern#1bp\raisebox{#2bp}[0pt][0pt]{\fontsize{#3bp}{#3bp}\csname font#4\endcsname\char#5}\hss}}
\newcommand{\RT}[4]{\hbox to 0pt{\kern#1bp\raisebox{#2bp}[0pt][0pt]{\fontsize{#3bp}{#3bp}\selectfont#4}\hss}}
\newcommand{\RS}[4]{\hbox to 0pt{\kern#1bp\raisebox{#2bp}[0pt][0pt]{\fontsize{#3bp}{#3bp}\selectfont\ensuremath{#4}}\hss}}
\newcommand{\RR}[4]{\hbox to 0pt{\kern#1bp\raisebox{#2bp}[0pt][0pt]{\rule{#3bp}{#4bp}}\hss}}
\newcommand{\fontmi}{\usefont{OML}{cmm}{m}{it}}
\newcommand{\fontsy}{\usefont{OMS}{cmsy}{m}{n}}
\newcommand{\fontex}{\usefont{OMX}{cmex}{m}{n}}
\newcommand{\fontrm}{\usefont{OT1}{cmr}{m}{n}}
\newcommand{\fontbf}{\usefont{OT1}{cmr}{bx}{n}}
\newcommand{\fontam}{\usefont{U}{msa}{m}{n}}
\newcommand{\fontbb}{\usefont{U}{msb}{m}{n}}
\widowpenalty=10000\clubpenalty=10000\raggedbottom
\begin{document}
\begin{center}{\LARGE\bfseries 22S09–22S15\par}\vspace{8pt}复旦大学高等代数习题集·题目与解答\end{center}
% Original page 1.
\section*{2022S09}
设\allowbreak{}\hskip.1em \RB{55.327}{8.685}{3.224}{%
\RG{0.000}{0.003}{10.9091}{mi}{86}\RG{11.818}{0.003}{10.9091}{rm}{61}\RG{23.333}{0.003}{10.9091}{mi}{77}\RG{33.916}{-1.634}{7.9701}{mi}{110}\RG{39.553}{0.003}{10.9091}{rm}{40}\RG{43.795}{0.003}{10.9091}{bb}{67}\RG{51.674}{0.003}{10.9091}{rm}{41}%
}\hskip.1em\allowbreak{} ，线性变换\allowbreak{}\hskip.1em \RB{49.662}{7.954}{2.875}{%
\RG{0.000}{0.003}{10.9091}{mi}{39}\RG{8.349}{0.003}{10.9091}{rm}{58}\RG{15.016}{0.003}{10.9091}{mi}{86}\RG{26.834}{0.003}{10.9091}{sy}{33}\RG{40.773}{0.003}{10.9091}{mi}{86}%
}\hskip.1em\allowbreak{} 定义为

\[
\RB{139.471}{8.681}{3.228}{%
\RG{0.000}{-0.001}{10.9091}{mi}{39}\RG{7.137}{-0.001}{10.9091}{rm}{40}\RG{11.379}{-0.001}{10.9091}{mi}{88}\RG{21.273}{-0.001}{10.9091}{rm}{41}\RG{28.528}{-0.001}{10.9091}{rm}{61}\RG{40.061}{-0.001}{10.9091}{mi}{74}\RG{47.152}{-0.001}{10.9091}{mi}{88}\RG{57.036}{-0.001}{10.9091}{mi}{74}\RG{62.927}{-0.001}{10.9091}{mi}{59}\RG{89.578}{-0.001}{10.9091}{mi}{74}\RG{99.736}{-0.001}{10.9091}{rm}{61}\RG{111.251}{-0.001}{10.9091}{mi}{74}\RG{122.937}{-0.001}{10.9091}{rm}{40}\RG{127.170}{-0.001}{10.9091}{rm}{48}\RG{132.624}{-0.001}{10.9091}{rm}{41}\RG{136.876}{-0.001}{10.9091}{mi}{58}\RG{117.300}{-1.638}{7.9701}{mi}{110}%
}
\]
其中\allowbreak{}\hskip.1em \RB{25.025}{8.687}{3.222}{%
\RG{0.000}{0.005}{10.9091}{mi}{74}\RG{6.049}{-1.631}{7.9701}{mi}{110}\RG{11.685}{0.005}{10.9091}{rm}{40}\RG{15.918}{0.005}{10.9091}{rm}{48}\RG{21.372}{0.005}{10.9091}{rm}{41}%
}\hskip.1em\allowbreak{} 为主对角线全为零、上副对角线全为\allowbreak{}\hskip.1em \RB{5.071}{7.771}{0.500}{%
\RG{0.000}{0.005}{10.9091}{rm}{49}%
}\hskip.1em\allowbreak{} 的Jordan块。求\allowbreak{}\hskip.1em \RB{7.253}{5.327}{2.873}{%
\RG{0.000}{0.005}{10.9091}{mi}{39}%
}\hskip.1em\allowbreak{} 的Jordan标准型。

解：记\allowbreak{}\hskip.1em \RB{14.722}{7.919}{3.761}{%
\RG{0.000}{0.001}{10.9091}{mi}{69}\RG{8.053}{-1.635}{7.9701}{mi}{105}\RG{10.930}{-1.635}{7.9701}{mi}{106}%
}\hskip.1em\allowbreak{} 为第\allowbreak{}\hskip.1em \RB{21.619}{8.683}{3.226}{%
\RG{0.000}{0.001}{10.9091}{rm}{40}\RG{4.242}{0.001}{10.9091}{mi}{105}\RG{7.995}{0.001}{10.9091}{mi}{59}\RG{12.828}{0.001}{10.9091}{mi}{106}\RG{17.966}{0.001}{10.9091}{rm}{41}%
}\hskip.1em\allowbreak{} 个元素等于\allowbreak{}\hskip.1em \RB{5.071}{7.767}{0.500}{%
\RG{0.000}{0.001}{10.9091}{rm}{49}%
}\hskip.1em\allowbreak{} 、其余元素等于零的基础矩阵。由矩阵乘法，

\[
\RB{271.754}{10.319}{2.125}{%
\RG{0.000}{1.637}{10.9091}{mi}{74}\RG{7.091}{1.637}{10.9091}{mi}{69}\RG{25.447}{1.637}{10.9091}{rm}{61}\RG{36.962}{1.637}{10.9091}{mi}{69}\RG{65.455}{1.637}{10.9091}{mi}{59}\RG{92.106}{1.637}{10.9091}{mi}{69}\RG{107.440}{1.637}{10.9091}{mi}{74}\RG{117.568}{1.637}{10.9091}{rm}{61}\RG{129.083}{1.637}{10.9091}{mi}{69}\RG{157.575}{1.637}{10.9091}{mi}{59}\RG{184.226}{1.637}{10.9091}{mi}{39}\RG{191.379}{1.637}{10.9091}{rm}{40}\RG{195.621}{1.637}{10.9091}{mi}{69}\RG{210.939}{1.637}{10.9091}{rm}{41}\RG{218.193}{1.637}{10.9091}{rm}{61}\RG{229.727}{1.637}{10.9091}{mi}{69}\RG{269.039}{1.637}{10.9091}{mi}{59}\RG{15.152}{0.001}{7.9701}{mi}{105}\RG{18.029}{0.001}{7.9701}{mi}{106}\RG{45.015}{0.001}{7.9701}{mi}{105}\RG{47.898}{0.001}{7.9701}{sy}{0}\RG{54.485}{0.001}{7.9701}{rm}{49}\RG{58.719}{0.001}{7.9701}{mi}{59}\RG{61.071}{0.001}{7.9701}{mi}{106}\RG{100.174}{0.001}{7.9701}{mi}{105}\RG{103.051}{0.001}{7.9701}{mi}{106}\RG{137.136}{0.001}{7.9701}{mi}{105}\RG{140.013}{0.001}{7.9701}{mi}{59}\RG{142.365}{0.001}{7.9701}{mi}{106}\RG{146.256}{0.001}{7.9701}{rm}{43}\RG{152.840}{0.001}{7.9701}{rm}{49}\RG{203.674}{0.001}{7.9701}{mi}{105}\RG{206.551}{0.001}{7.9701}{mi}{106}\RG{237.780}{0.001}{7.9701}{mi}{105}\RG{240.663}{0.001}{7.9701}{sy}{0}\RG{247.250}{0.001}{7.9701}{rm}{49}\RG{251.484}{0.001}{7.9701}{mi}{59}\RG{253.835}{0.001}{7.9701}{mi}{106}\RG{257.721}{0.001}{7.9701}{rm}{43}\RG{264.305}{0.001}{7.9701}{rm}{49}%
}
\]
其中越界的基础矩阵约定为零。因此，\allowbreak{}\hskip.1em \RB{7.253}{5.325}{2.875}{%
\RG{0.000}{0.003}{10.9091}{mi}{39}%
}\hskip.1em\allowbreak{} 将每条反对角线\allowbreak{}\hskip.1em \RB{42.052}{7.714}{2.733}{%
\RG{0.000}{0.003}{10.9091}{mi}{105}\RG{6.214}{0.003}{10.9091}{rm}{43}\RG{17.154}{0.003}{10.9091}{mi}{106}\RG{25.378}{0.003}{10.9091}{rm}{61}\RG{36.970}{0.003}{10.9091}{mi}{115}%
}\hskip.1em\allowbreak{} 上的矩阵仍映到同一条反对角线上。

先取\allowbreak{}\hskip.1em \RB{56.151}{7.769}{2.602}{%
\RG{0.000}{0.004}{10.9091}{mi}{114}\RG{8.255}{0.004}{10.9091}{rm}{61}\RG{19.753}{0.004}{10.9091}{rm}{49}\RG{25.225}{0.004}{10.9091}{mi}{59}\RG{30.058}{0.004}{10.9091}{mi}{58}\RG{34.902}{0.004}{10.9091}{mi}{58}\RG{39.734}{0.004}{10.9091}{mi}{58}\RG{44.578}{0.004}{10.9091}{mi}{59}\RG{49.422}{0.004}{10.9091}{mi}{110}%
}\hskip.1em\allowbreak{} ，得到长度为\allowbreak{}\hskip.1em \RB{5.256}{5.326}{0.616}{%
\RG{0.000}{0.004}{10.9091}{mi}{114}%
}\hskip.1em\allowbreak{} 的链

\[
\RB{183.256}{7.918}{3.674}{%
\RG{0.000}{-0.000}{10.9091}{mi}{69}\RG{21.754}{-0.000}{10.9091}{sy}{55}\RG{21.754}{-0.000}{10.9091}{sy}{0}\RG{28.420}{-0.000}{10.9091}{sy}{33}\RG{42.360}{-0.000}{10.9091}{mi}{69}\RG{75.406}{-0.000}{10.9091}{sy}{55}\RG{75.406}{-0.000}{10.9091}{sy}{0}\RG{82.071}{-0.000}{10.9091}{sy}{33}\RG{96.002}{-0.000}{10.9091}{sy}{1}\RG{100.835}{-0.000}{10.9091}{sy}{1}\RG{105.679}{-0.000}{10.9091}{sy}{1}\RG{111.733}{-0.000}{10.9091}{sy}{55}\RG{111.733}{-0.000}{10.9091}{sy}{0}\RG{118.399}{-0.000}{10.9091}{sy}{33}\RG{132.376}{-0.000}{10.9091}{mi}{69}\RG{154.600}{-0.000}{10.9091}{sy}{55}\RG{154.600}{-0.000}{10.9091}{sy}{0}\RG{161.266}{-0.000}{10.9091}{sy}{33}\RG{175.206}{-0.000}{10.9091}{rm}{48}\RG{180.661}{-0.000}{10.9091}{mi}{58}\RG{8.053}{-1.636}{7.9701}{mi}{114}\RG{11.640}{-1.636}{7.9701}{mi}{59}\RG{13.991}{-1.636}{7.9701}{rm}{49}\RG{50.413}{-1.636}{7.9701}{mi}{114}\RG{54.469}{-1.636}{7.9701}{sy}{0}\RG{61.057}{-1.636}{7.9701}{rm}{49}\RG{65.291}{-1.636}{7.9701}{mi}{59}\RG{67.643}{-1.636}{7.9701}{rm}{50}\RG{140.429}{-1.636}{7.9701}{rm}{49}\RG{144.663}{-1.636}{7.9701}{mi}{59}\RG{147.014}{-1.636}{7.9701}{mi}{114}%
}
\]
再取\allowbreak{}\hskip.1em \RB{54.865}{7.763}{2.608}{%
\RG{0.000}{-0.003}{10.9091}{mi}{116}\RG{6.969}{-0.003}{10.9091}{rm}{61}\RG{18.467}{-0.003}{10.9091}{rm}{50}\RG{23.939}{-0.003}{10.9091}{mi}{59}\RG{28.772}{-0.003}{10.9091}{mi}{58}\RG{33.616}{-0.003}{10.9091}{mi}{58}\RG{38.448}{-0.003}{10.9091}{mi}{58}\RG{43.292}{-0.003}{10.9091}{mi}{59}\RG{48.136}{-0.003}{10.9091}{mi}{110}%
}\hskip.1em\allowbreak{} ，得到长度为\allowbreak{}\hskip.1em \RB{42.214}{7.763}{1.408}{%
\RG{0.000}{-0.003}{10.9091}{mi}{110}\RG{8.972}{-0.003}{10.9091}{sy}{0}\RG{19.881}{-0.003}{10.9091}{mi}{116}\RG{26.245}{-0.003}{10.9091}{rm}{43}\RG{37.143}{-0.003}{10.9091}{rm}{49}%
}\hskip.1em\allowbreak{} 的链

\[
\RB{194.263}{7.921}{3.672}{%
\RG{0.000}{0.003}{10.9091}{mi}{69}\RG{22.130}{0.003}{10.9091}{sy}{55}\RG{22.130}{0.003}{10.9091}{sy}{0}\RG{28.796}{0.003}{10.9091}{sy}{33}\RG{42.736}{0.003}{10.9091}{mi}{69}\RG{86.508}{0.003}{10.9091}{sy}{55}\RG{86.508}{0.003}{10.9091}{sy}{0}\RG{93.173}{0.003}{10.9091}{sy}{33}\RG{107.104}{0.003}{10.9091}{sy}{1}\RG{111.937}{0.003}{10.9091}{sy}{1}\RG{116.781}{0.003}{10.9091}{sy}{1}\RG{122.835}{0.003}{10.9091}{sy}{55}\RG{122.835}{0.003}{10.9091}{sy}{0}\RG{129.501}{0.003}{10.9091}{sy}{33}\RG{143.478}{0.003}{10.9091}{mi}{69}\RG{165.608}{0.003}{10.9091}{sy}{55}\RG{165.608}{0.003}{10.9091}{sy}{0}\RG{172.274}{0.003}{10.9091}{sy}{33}\RG{186.215}{0.003}{10.9091}{rm}{48}\RG{191.669}{0.003}{10.9091}{mi}{58}\RG{8.053}{-1.634}{7.9701}{mi}{110}\RG{13.186}{-1.634}{7.9701}{mi}{59}\RG{15.537}{-1.634}{7.9701}{mi}{116}\RG{50.789}{-1.634}{7.9701}{mi}{110}\RG{55.927}{-1.634}{7.9701}{sy}{0}\RG{62.514}{-1.634}{7.9701}{rm}{49}\RG{66.749}{-1.634}{7.9701}{mi}{59}\RG{69.100}{-1.634}{7.9701}{mi}{116}\RG{72.159}{-1.634}{7.9701}{rm}{43}\RG{78.742}{-1.634}{7.9701}{rm}{49}\RG{151.531}{-1.634}{7.9701}{mi}{116}\RG{154.584}{-1.634}{7.9701}{mi}{59}\RG{156.935}{-1.634}{7.9701}{mi}{110}%
}
\]
第一组链对应\allowbreak{}\hskip.1em \RB{90.058}{7.765}{2.737}{%
\RG{0.000}{-0.001}{10.9091}{mi}{105}\RG{5.850}{-0.001}{10.9091}{rm}{43}\RG{16.427}{-0.001}{10.9091}{mi}{106}\RG{24.574}{-0.001}{10.9091}{rm}{61}\RG{36.073}{-0.001}{10.9091}{rm}{50}\RG{41.544}{-0.001}{10.9091}{mi}{59}\RG{46.377}{-0.001}{10.9091}{mi}{58}\RG{51.221}{-0.001}{10.9091}{mi}{58}\RG{56.064}{-0.001}{10.9091}{mi}{58}\RG{60.897}{-0.001}{10.9091}{mi}{59}\RG{65.741}{-0.001}{10.9091}{mi}{110}\RG{74.427}{-0.001}{10.9091}{rm}{43}\RG{84.987}{-0.001}{10.9091}{rm}{49}%
}\hskip.1em\allowbreak{} ，第二组链对应\allowbreak{}\hskip.1em \RB{97.188}{7.765}{2.737}{%
\RG{0.000}{-0.001}{10.9091}{mi}{105}\RG{5.850}{-0.001}{10.9091}{rm}{43}\RG{16.427}{-0.001}{10.9091}{mi}{106}\RG{24.574}{-0.001}{10.9091}{rm}{61}\RG{36.089}{-0.001}{10.9091}{mi}{110}\RG{44.729}{-0.001}{10.9091}{rm}{43}\RG{55.290}{-0.001}{10.9091}{rm}{50}\RG{60.762}{-0.001}{10.9091}{mi}{59}\RG{65.594}{-0.001}{10.9091}{mi}{58}\RG{70.438}{-0.001}{10.9091}{mi}{58}\RG{75.271}{-0.001}{10.9091}{mi}{58}\RG{80.114}{-0.001}{10.9091}{mi}{59}\RG{85.004}{-0.001}{10.9091}{rm}{50}\RG{90.459}{-0.001}{10.9091}{mi}{110}%
}\hskip.1em\allowbreak{} 。故这些链互不重叠，并且恰好包含全部\allowbreak{}\hskip.1em \RB{10.850}{9.762}{0.616}{%
\RG{0.000}{0.004}{10.9091}{mi}{110}\RG{6.548}{3.962}{7.9701}{rm}{50}%
}\hskip.1em\allowbreak{} 个基础矩阵。

将每条链的非零向量按上述箭头的逆序排列，则在所得基下，每条长度为\allowbreak{}\hskip.1em \RB{6.042}{8.071}{0.620}{%
\RG{0.000}{0.000}{10.9091}{mi}{107}%
}\hskip.1em\allowbreak{} 的链对应一个\allowbreak{}\hskip.1em \RB{24.509}{8.677}{3.232}{%
\RG{0.000}{-0.005}{10.9091}{mi}{74}\RG{6.049}{-1.782}{7.9701}{mi}{107}\RG{11.169}{-0.005}{10.9091}{rm}{40}\RG{15.402}{-0.005}{10.9091}{rm}{48}\RG{20.856}{-0.005}{10.9091}{rm}{41}%
}\hskip.1em\allowbreak{} 。因此所求标准型为

\[
\RB{138.855}{19.434}{13.935}{%
\RG{46.487}{13.634}{7.9701}{mi}{110}\RG{51.625}{13.634}{7.9701}{sy}{0}\RG{58.213}{13.634}{7.9701}{rm}{49}\RG{46.224}{10.361}{10.9092}{ex}{77}\RG{62.710}{8.834}{10.9092}{ex}{0}\RG{131.260}{8.834}{10.9092}{ex}{1}\RG{0.000}{-0.002}{10.9091}{mi}{74}\RG{11.686}{-0.002}{10.9091}{rm}{40}\RG{15.919}{-0.002}{10.9091}{rm}{48}\RG{21.373}{-0.002}{10.9091}{rm}{41}\RG{31.682}{-0.002}{10.9091}{sy}{8}\RG{67.710}{-0.002}{10.9091}{mi}{74}\RG{78.879}{-0.002}{10.9091}{rm}{40}\RG{83.111}{-0.002}{10.9091}{rm}{48}\RG{88.566}{-0.002}{10.9091}{rm}{41}\RG{95.243}{-0.002}{10.9091}{sy}{8}\RG{106.152}{-0.002}{10.9091}{mi}{74}\RG{117.321}{-0.002}{10.9091}{rm}{40}\RG{121.554}{-0.002}{10.9091}{rm}{48}\RG{127.008}{-0.002}{10.9091}{rm}{41}\RG{136.260}{-0.002}{10.9091}{mi}{58}\RG{6.050}{-1.639}{7.9701}{mi}{110}\RG{73.759}{-1.780}{7.9701}{mi}{107}\RG{112.201}{-1.780}{7.9701}{mi}{107}\RG{46.745}{-13.355}{7.9701}{mi}{107}\RG{51.367}{-13.355}{7.9701}{rm}{61}\RG{57.951}{-13.355}{7.9701}{rm}{49}%
}
\]
即\allowbreak{}\hskip.1em \RB{6.729}{5.325}{0.617}{%
\RG{0.000}{0.003}{10.9091}{mi}{110}%
}\hskip.1em\allowbreak{} 阶块恰有一个，而\allowbreak{}\hskip.1em \RB{64.952}{7.769}{2.602}{%
\RG{0.000}{0.003}{10.9091}{rm}{49}\RG{5.454}{0.003}{10.9091}{mi}{59}\RG{10.303}{0.003}{10.9091}{rm}{50}\RG{15.757}{0.003}{10.9091}{mi}{59}\RG{20.590}{0.003}{10.9091}{mi}{58}\RG{25.434}{0.003}{10.9091}{mi}{58}\RG{30.277}{0.003}{10.9091}{mi}{58}\RG{35.110}{0.003}{10.9091}{mi}{59}\RG{39.954}{0.003}{10.9091}{mi}{110}\RG{48.972}{0.003}{10.9091}{sy}{0}\RG{59.882}{0.003}{10.9091}{rm}{49}%
}\hskip.1em\allowbreak{} 阶块各有两个。维数核对为

\[
\RB{79.374}{19.439}{13.930}{%
\RG{27.154}{13.639}{7.9701}{mi}{110}\RG{32.292}{13.639}{7.9701}{sy}{0}\RG{38.879}{13.639}{7.9701}{rm}{49}\RG{27.255}{10.366}{10.9092}{ex}{88}\RG{72.048}{4.507}{7.9701}{rm}{50}\RG{0.000}{0.003}{10.9091}{mi}{110}\RG{8.972}{0.003}{10.9091}{rm}{43}\RG{19.870}{0.003}{10.9091}{rm}{50}\RG{44.931}{0.003}{10.9091}{mi}{107}\RG{53.984}{0.003}{10.9091}{rm}{61}\RG{65.500}{0.003}{10.9091}{mi}{110}\RG{76.780}{0.003}{10.9091}{mi}{58}\RG{27.412}{-13.350}{7.9701}{mi}{107}\RG{32.034}{-13.350}{7.9701}{rm}{61}\RG{38.618}{-13.350}{7.9701}{rm}{49}%
}
\]
还可用幂的秩作独立核对。对\allowbreak{}\hskip.1em \RB{47.297}{8.068}{1.998}{%
\RG{0.000}{-0.003}{10.9091}{rm}{48}\RG{8.485}{-0.003}{10.9091}{sy}{20}\RG{20.000}{-0.003}{10.9091}{mi}{107}\RG{29.053}{-0.003}{10.9091}{sy}{20}\RG{40.568}{-0.003}{10.9091}{mi}{110}%
}\hskip.1em\allowbreak{} ，

\[
\RB{93.286}{10.538}{3.901}{%
\RG{7.136}{4.507}{7.9701}{mi}{107}\RG{0.000}{0.003}{10.9091}{mi}{39}\RG{12.256}{0.003}{10.9091}{rm}{40}\RG{16.498}{0.003}{10.9091}{mi}{69}\RG{31.817}{0.003}{10.9091}{rm}{41}\RG{39.072}{0.003}{10.9091}{rm}{61}\RG{50.605}{0.003}{10.9091}{mi}{69}\RG{90.692}{0.003}{10.9091}{mi}{58}\RG{24.551}{-1.634}{7.9701}{mi}{105}\RG{27.428}{-1.634}{7.9701}{mi}{106}\RG{58.658}{-1.775}{7.9701}{mi}{105}\RG{61.541}{-1.775}{7.9701}{sy}{0}\RG{68.128}{-1.775}{7.9701}{mi}{107}\RG{72.742}{-1.775}{7.9701}{mi}{59}\RG{75.093}{-1.775}{7.9701}{mi}{106}\RG{78.986}{-1.775}{7.9701}{rm}{43}\RG{85.573}{-1.775}{7.9701}{mi}{107}%
}
\]
像非零当且仅当\allowbreak{}\hskip.1em \RB{24.326}{8.070}{0.937}{%
\RG{0.000}{-0.001}{10.9091}{mi}{105}\RG{6.775}{-0.001}{10.9091}{mi}{62}\RG{18.284}{-0.001}{10.9091}{mi}{107}%
}\hskip.1em\allowbreak{} 且\allowbreak{}\hskip.1em \RB{45.586}{8.070}{2.737}{%
\RG{0.000}{-0.001}{10.9091}{mi}{106}\RG{8.147}{-0.001}{10.9091}{sy}{20}\RG{19.662}{-0.001}{10.9091}{mi}{110}\RG{28.635}{-0.001}{10.9091}{sy}{0}\RG{39.544}{-0.001}{10.9091}{mi}{107}%
}\hskip.1em\allowbreak{} ；这些非零像彼此不同，故

\[
\RB{98.559}{10.540}{3.222}{%
\RG{32.935}{4.509}{7.9701}{mi}{107}\RG{91.232}{4.509}{7.9701}{rm}{50}\RT{0.000}{0.005}{10.9091}{r}\RT{4.276}{0.005}{10.9091}{a}\RT{9.731}{0.005}{10.9091}{n}\RT{15.797}{0.005}{10.9091}{k}\RG{21.556}{0.005}{10.9091}{rm}{40}\RG{25.798}{0.005}{10.9091}{mi}{39}\RG{38.055}{0.005}{10.9091}{rm}{41}\RG{45.310}{0.005}{10.9091}{rm}{61}\RG{56.819}{0.005}{10.9091}{rm}{40}\RG{61.086}{0.005}{10.9091}{mi}{110}\RG{70.058}{0.005}{10.9091}{sy}{0}\RG{80.967}{0.005}{10.9091}{mi}{107}\RG{86.990}{0.005}{10.9091}{rm}{41}\RG{95.965}{0.005}{10.9091}{mi}{58}%
}
\]
这也与上述Jordan块完全一致。当\allowbreak{}\hskip.1em \RB{26.147}{7.767}{0.619}{%
\RG{0.000}{0.001}{10.9091}{mi}{110}\RG{9.578}{0.001}{10.9091}{rm}{61}\RG{21.076}{0.001}{10.9091}{rm}{49}%
}\hskip.1em\allowbreak{} 时，直和中的求和部分为空，结果就是\allowbreak{}\hskip.1em \RB{24.121}{8.683}{3.226}{%
\RG{0.000}{0.001}{10.9091}{mi}{74}\RG{6.049}{-1.636}{7.9701}{rm}{49}\RG{10.781}{0.001}{10.9091}{rm}{40}\RG{15.014}{0.001}{10.9091}{rm}{48}\RG{20.468}{0.001}{10.9091}{rm}{41}%
}\hskip.1em\allowbreak{} 。\allowbreak{}\hskip.1em \RB{8.366}{7.996}{0.500}{%
\RG{0.000}{0.001}{10.9091}{am}{3}%
}\hskip.1em\allowbreak{}

方法小结：先考察线性变换对基础矩阵的作用，再按反对角线组织Jordan链；链的长度直接给出Jordan块的阶数与个数。

题目来源：谢启鸿《复旦大学高等代数习题集（2023版）》，第116–117页。

% Original page 2.
\section*{2022S10}
设\allowbreak{}\hskip.1em \RB{27.268}{7.976}{0.933}{%
\RG{0.000}{0.003}{10.9091}{mi}{97}\RG{8.796}{0.003}{10.9091}{sy}{50}\RG{19.099}{0.003}{10.9091}{bb}{82}%
}\hskip.1em\allowbreak{} ，求实对称矩阵

\[
\RB{182.423}{71.425}{8.099}{%
\RG{22.727}{70.500}{10.9092}{ex}{48}\RG{175.563}{70.500}{10.9092}{ex}{49}\RG{108.275}{65.178}{7.9701}{rm}{50}\RG{159.106}{65.178}{7.9701}{mi}{110}\RG{164.245}{65.178}{7.9701}{sy}{0}\RG{170.831}{65.178}{7.9701}{rm}{49}\RG{40.657}{61.220}{10.9091}{rm}{49}\RG{72.688}{61.220}{10.9091}{mi}{97}\RG{102.492}{61.220}{10.9091}{mi}{97}\RG{128.832}{61.220}{10.9091}{sy}{1}\RG{133.665}{61.220}{10.9091}{sy}{1}\RG{138.509}{61.220}{10.9091}{sy}{1}\RG{153.340}{61.220}{10.9091}{mi}{97}\RG{22.727}{51.299}{10.9092}{ex}{66}\RG{175.563}{51.299}{10.9092}{ex}{67}\RG{159.106}{50.002}{7.9701}{mi}{110}\RG{164.245}{50.002}{7.9701}{sy}{0}\RG{170.831}{50.002}{7.9701}{rm}{50}\RG{40.501}{46.044}{10.9091}{mi}{97}\RG{72.844}{46.044}{10.9091}{rm}{49}\RG{104.874}{46.044}{10.9091}{mi}{97}\RG{128.832}{46.044}{10.9091}{sy}{1}\RG{133.665}{46.044}{10.9091}{sy}{1}\RG{138.509}{46.044}{10.9091}{sy}{1}\RG{153.340}{46.044}{10.9091}{mi}{97}\RG{22.727}{44.754}{10.9092}{ex}{66}\RG{175.563}{44.754}{10.9092}{ex}{67}\RG{22.727}{38.208}{10.9092}{ex}{66}\RG{175.563}{38.208}{10.9092}{ex}{67}\RG{43.901}{34.827}{7.9701}{rm}{50}\RG{159.106}{34.827}{7.9701}{mi}{110}\RG{164.245}{34.827}{7.9701}{sy}{0}\RG{170.831}{34.827}{7.9701}{rm}{51}\RG{22.727}{31.663}{10.9092}{ex}{66}\RG{175.563}{31.663}{10.9092}{ex}{67}\RG{38.135}{30.869}{10.9091}{mi}{97}\RG{72.688}{30.869}{10.9091}{mi}{97}\RG{105.030}{30.869}{10.9091}{rm}{49}\RG{128.832}{30.869}{10.9091}{sy}{1}\RG{133.665}{30.869}{10.9091}{sy}{1}\RG{138.508}{30.869}{10.9091}{sy}{1}\RG{153.340}{30.869}{10.9091}{mi}{97}\RG{0.000}{28.935}{10.9091}{mi}{65}\RG{11.212}{28.935}{10.9091}{rm}{61}\RG{22.727}{25.117}{10.9092}{ex}{66}\RG{175.563}{25.117}{10.9092}{ex}{67}\RT{41.868}{19.182}{10.9091}{.}\RT{74.055}{19.182}{10.9091}{.}\RT{106.241}{19.182}{10.9091}{.}\RT{162.935}{19.182}{10.9091}{.}\RG{22.727}{18.572}{10.9092}{ex}{66}\RG{175.563}{18.572}{10.9092}{ex}{67}\RT{130.343}{18.186}{10.9091}{.}\RT{41.868}{15.197}{10.9091}{.}\RT{74.055}{15.197}{10.9091}{.}\RT{106.241}{15.197}{10.9091}{.}\RT{134.588}{15.197}{10.9091}{.}\RT{162.935}{15.197}{10.9091}{.}\RT{138.833}{12.208}{10.9091}{.}\RG{22.727}{11.590}{10.9092}{ex}{64}\RG{175.563}{11.590}{10.9092}{ex}{65}\RT{41.868}{11.212}{10.9091}{.}\RT{74.055}{11.212}{10.9091}{.}\RT{106.241}{11.212}{10.9091}{.}\RT{162.935}{11.212}{10.9091}{.}\RG{38.039}{-0.004}{7.9701}{mi}{110}\RG{43.177}{-0.004}{7.9701}{sy}{0}\RG{49.764}{-0.004}{7.9701}{rm}{49}\RG{70.226}{-0.004}{7.9701}{mi}{110}\RG{75.364}{-0.004}{7.9701}{sy}{0}\RG{81.951}{-0.004}{7.9701}{rm}{50}\RG{102.412}{-0.004}{7.9701}{mi}{110}\RG{107.550}{-0.004}{7.9701}{sy}{0}\RG{114.137}{-0.004}{7.9701}{rm}{51}\RG{32.273}{-3.963}{10.9091}{mi}{97}\RG{64.460}{-3.963}{10.9091}{mi}{97}\RG{96.646}{-3.963}{10.9091}{mi}{97}\RG{128.832}{-3.963}{10.9091}{sy}{1}\RG{133.665}{-3.963}{10.9091}{sy}{1}\RG{138.508}{-3.963}{10.9091}{sy}{1}\RG{161.724}{-3.963}{10.9091}{rm}{49}%
}
\]
的正、负惯性指数。解：记\allowbreak{}\hskip.1em \RB{87.439}{10.436}{3.763}{%
\RG{0.000}{-0.001}{10.9091}{mi}{65}\RG{11.212}{-0.001}{10.9091}{rm}{61}\RG{22.710}{-0.001}{10.9091}{rm}{40}\RG{26.969}{-0.001}{10.9091}{mi}{97}\RG{32.736}{3.958}{7.9701}{sy}{106}\RG{35.088}{3.958}{7.9701}{mi}{105}\RG{37.971}{3.958}{7.9701}{sy}{0}\RG{44.558}{3.958}{7.9701}{mi}{106}\RG{48.442}{3.958}{7.9701}{sy}{106}\RG{51.292}{-0.001}{10.9091}{rm}{41}\RG{55.535}{-1.637}{7.9701}{rm}{49}\RG{59.769}{-1.637}{7.9701}{sy}{20}\RG{66.356}{-1.637}{7.9701}{mi}{105}\RG{69.233}{-1.637}{7.9701}{mi}{59}\RG{71.584}{-1.637}{7.9701}{mi}{106}\RG{75.476}{-1.637}{7.9701}{sy}{20}\RG{82.062}{-1.637}{7.9701}{mi}{110}%
}\hskip.1em\allowbreak{} 。取可逆下三角矩阵

\[
\RB{221.151}{42.493}{37.032}{%
\RG{23.880}{41.567}{10.9092}{ex}{48}\RG{134.220}{41.567}{10.9092}{ex}{49}\RG{37.824}{32.286}{10.9091}{rm}{49}\RG{23.880}{22.366}{10.9092}{ex}{66}\RG{134.220}{22.366}{10.9092}{ex}{67}\RG{33.426}{17.111}{10.9091}{sy}{0}\RG{41.911}{17.111}{10.9091}{mi}{97}\RG{62.039}{17.111}{10.9091}{rm}{49}\RG{23.880}{15.821}{10.9092}{ex}{66}\RG{134.220}{15.821}{10.9092}{ex}{67}\RG{23.880}{9.275}{10.9092}{ex}{66}\RG{134.220}{9.275}{10.9092}{ex}{67}\RG{23.880}{2.730}{10.9092}{ex}{66}\RG{134.220}{2.730}{10.9092}{ex}{67}\RG{57.640}{1.936}{10.9091}{sy}{0}\RG{66.126}{1.936}{10.9091}{mi}{97}\RG{85.495}{1.936}{10.9091}{rm}{49}\RG{0.000}{0.003}{10.9091}{mi}{68}\RG{12.365}{0.003}{10.9091}{rm}{61}\RG{145.584}{0.003}{10.9091}{mi}{59}\RT{172.251}{0.003}{10.9091}{d}\RT{178.317}{0.003}{10.9091}{e}\RT{183.161}{0.003}{10.9091}{t}\RG{189.222}{0.003}{10.9091}{mi}{68}\RG{201.587}{0.003}{10.9091}{rm}{61}\RG{213.085}{0.003}{10.9091}{rm}{49}\RG{218.557}{0.003}{10.9091}{mi}{58}\RG{23.880}{-3.816}{10.9092}{ex}{66}\RG{134.220}{-3.816}{10.9092}{ex}{67}\RG{23.880}{-10.361}{10.9092}{ex}{66}\RG{134.220}{-10.361}{10.9092}{ex}{67}\RT{82.461}{-10.747}{10.9091}{.}\RT{105.916}{-10.747}{10.9091}{.}\RT{86.706}{-13.736}{10.9091}{.}\RT{110.161}{-13.736}{10.9091}{.}\RT{90.950}{-16.725}{10.9091}{.}\RT{114.406}{-16.725}{10.9091}{.}\RG{23.880}{-17.343}{10.9092}{ex}{64}\RG{134.220}{-17.343}{10.9092}{ex}{65}\RG{104.552}{-32.896}{10.9091}{sy}{0}\RG{113.037}{-32.896}{10.9091}{mi}{97}\RG{128.766}{-32.896}{10.9091}{rm}{49}%
}
\]
左乘\allowbreak{}\hskip.1em \RB{9.271}{7.952}{0.500}{%
\RG{0.000}{0.001}{10.9091}{mi}{68}%
}\hskip.1em\allowbreak{} 相当于用原矩阵的第\allowbreak{}\hskip.1em \RB{3.696}{7.712}{0.619}{%
\RG{0.000}{0.001}{10.9091}{mi}{105}%
}\hskip.1em\allowbreak{} 行减去第\allowbreak{}\hskip.1em \RB{22.245}{7.766}{0.619}{%
\RG{0.000}{0.001}{10.9091}{mi}{105}\RG{6.224}{0.001}{10.9091}{sy}{0}\RG{17.174}{0.001}{10.9091}{rm}{49}%
}\hskip.1em\allowbreak{} 行的\allowbreak{}\hskip.1em \RB{5.933}{5.323}{0.619}{%
\RG{0.000}{0.001}{10.9091}{mi}{97}%
}\hskip.1em\allowbreak{} 倍，\allowbreak{}\hskip.1em \RB{54.889}{7.766}{2.605}{%
\RG{0.000}{0.001}{10.9091}{mi}{105}\RG{6.891}{0.001}{10.9091}{rm}{61}\RG{18.498}{0.001}{10.9091}{rm}{50}\RG{23.963}{0.001}{10.9091}{mi}{59}\RG{28.796}{0.001}{10.9091}{mi}{58}\RG{33.640}{0.001}{10.9091}{mi}{58}\RG{38.483}{0.001}{10.9091}{mi}{58}\RG{43.316}{0.001}{10.9091}{mi}{59}\RG{48.160}{0.001}{10.9091}{mi}{110}%
}\hskip.1em\allowbreak{} 。具体地，第一行不变；对

\[
\RB{309.143}{44.714}{7.378}{%
\RG{0.000}{35.486}{10.9091}{mi}{105}\RG{6.788}{35.486}{10.9091}{sy}{21}\RG{18.304}{35.486}{10.9091}{rm}{50}\RT{23.758}{35.486}{10.9091}{，}\RG{202.013}{32.384}{10.9092}{ex}{56}\RG{202.013}{22.566}{10.9092}{ex}{60}\RG{211.710}{18.211}{10.9091}{rm}{48}\RG{217.164}{18.211}{10.9091}{mi}{59}\RG{283.039}{18.211}{10.9091}{mi}{106}\RG{291.167}{18.211}{10.9091}{mi}{60}\RG{302.676}{18.211}{10.9091}{mi}{105}\RG{306.429}{18.211}{10.9091}{mi}{59}\RG{154.199}{10.020}{10.9091}{rm}{40}\RG{158.442}{10.020}{10.9091}{mi}{68}\RG{167.759}{10.020}{10.9091}{mi}{65}\RG{175.959}{10.020}{10.9091}{rm}{41}\RG{190.498}{10.020}{10.9091}{rm}{61}\RG{180.201}{8.384}{7.9701}{mi}{105}\RG{183.079}{8.384}{7.9701}{mi}{106}\RG{240.506}{3.960}{7.9701}{rm}{50}\RG{255.248}{3.960}{7.9701}{mi}{106}\RG{259.132}{3.960}{7.9701}{sy}{0}\RG{265.718}{3.960}{7.9701}{mi}{105}\RG{202.013}{2.929}{10.9092}{ex}{58}\RG{211.710}{0.001}{10.9091}{rm}{40}\RG{215.943}{0.001}{10.9091}{rm}{49}\RG{223.831}{0.001}{10.9091}{sy}{0}\RG{234.740}{0.001}{10.9091}{mi}{97}\RG{245.239}{0.001}{10.9091}{rm}{41}\RG{249.481}{0.001}{10.9091}{mi}{97}\RG{269.099}{0.001}{10.9091}{mi}{59}\RG{283.030}{0.001}{10.9091}{mi}{106}\RG{291.186}{0.001}{10.9091}{sy}{21}\RG{302.702}{0.001}{10.9091}{mi}{105}\RG{306.454}{0.001}{10.9091}{mi}{58}%
}
\]
再右乘\allowbreak{}\hskip.1em \RB{9.271}{7.946}{0.505}{%
\RG{0.000}{-0.005}{10.9091}{mi}{68}%
}\hskip.1em\allowbreak{} T，即进行相同的列变换。对\allowbreak{}\hskip.1em \RB{25.060}{7.761}{2.741}{%
\RG{0.000}{-0.005}{10.9091}{mi}{106}\RG{8.147}{-0.005}{10.9091}{sy}{21}\RG{19.662}{-0.005}{10.9091}{rm}{50}%
}\hskip.1em\allowbreak{} ，第\allowbreak{}\hskip.1em \RB{4.831}{7.706}{2.741}{%
\RG{0.000}{-0.005}{10.9091}{mi}{106}%
}\hskip.1em\allowbreak{} 列减去原第\allowbreak{}\hskip.1em \RB{23.522}{7.761}{2.741}{%
\RG{0.000}{-0.005}{10.9091}{mi}{106}\RG{7.541}{-0.005}{10.9091}{sy}{0}\RG{18.451}{-0.005}{10.9091}{rm}{49}%
}\hskip.1em\allowbreak{} 列的\allowbreak{}\hskip.1em \RB{5.933}{5.317}{0.625}{%
\RG{0.000}{-0.005}{10.9091}{mi}{97}%
}\hskip.1em\allowbreak{} 倍，因而

\[
\RB{171.888}{11.429}{3.226}{%
\RT{26.851}{4.505}{7.9701}{T}\RG{106.789}{4.505}{7.9701}{rm}{50}\RG{160.318}{4.505}{7.9701}{rm}{50}\RG{0.000}{0.001}{10.9091}{mi}{68}\RG{9.316}{0.001}{10.9091}{mi}{65}\RG{17.498}{0.001}{10.9091}{mi}{68}\RG{36.166}{0.001}{10.9091}{rm}{61}\RT{47.681}{0.001}{10.9091}{d}\RT{53.747}{0.001}{10.9091}{i}\RT{56.780}{0.001}{10.9091}{a}\RT{62.234}{0.001}{10.9091}{g}\RG{67.689}{0.001}{10.9091}{rm}{40}\RG{71.921}{0.001}{10.9091}{rm}{49}\RG{77.386}{0.001}{10.9091}{mi}{59}\RG{82.235}{0.001}{10.9091}{rm}{49}\RG{90.113}{0.001}{10.9091}{sy}{0}\RG{101.023}{0.001}{10.9091}{mi}{97}\RG{111.521}{0.001}{10.9091}{mi}{59}\RG{116.354}{0.001}{10.9091}{mi}{58}\RG{121.197}{0.001}{10.9091}{mi}{58}\RG{126.041}{0.001}{10.9091}{mi}{58}\RG{130.874}{0.001}{10.9091}{mi}{59}\RG{135.764}{0.001}{10.9091}{rm}{49}\RG{143.642}{0.001}{10.9091}{sy}{0}\RG{154.551}{0.001}{10.9091}{mi}{97}\RG{165.050}{0.001}{10.9091}{rm}{41}\RG{169.293}{0.001}{10.9091}{mi}{58}%
}
\]
这里\allowbreak{}\hskip.1em \RB{26.147}{7.763}{0.623}{%
\RG{0.000}{-0.003}{10.9091}{mi}{110}\RG{9.578}{-0.003}{10.9091}{rm}{61}\RG{21.076}{-0.003}{10.9091}{rm}{49}%
}\hskip.1em\allowbreak{} 时，右端只有第一个对角元\allowbreak{}\hskip.1em \RB{5.071}{7.763}{0.503}{%
\RG{0.000}{-0.003}{10.9091}{rm}{49}%
}\hskip.1em\allowbreak{} 。

合同变换对应二次型的可逆变量替换，保持其正定子空间、负定子空间的最大维数，因而保持正、负惯性指数。这正是惯性定律的含义。因此，若同时记零惯性指数为\allowbreak{}\hskip.1em \RB{10.945}{5.320}{2.305}{%
\RG{0.000}{-0.002}{10.9091}{mi}{110}\RG{6.548}{-1.638}{7.9701}{rm}{48}%
}\hskip.1em\allowbreak{} ，则

参\allowbreak{}\hskip.1em \RB{8.964}{0.500}{18.805}{%
\RG{0.000}{-15.578}{10.9091}{sy}{106}\RG{3.031}{-15.578}{10.9091}{mi}{97}%
}\hskip.1em\allowbreak{} 数\allowbreak{}\hskip.1em \RB{14.131}{0.500}{18.805}{%
\RG{0.000}{-15.578}{10.9091}{sy}{106}\RG{6.060}{-15.578}{10.9091}{mi}{60}%
}\hskip.1em\allowbreak{} 条\allowbreak{}\hskip.1em \RB{5.071}{0.500}{16.078}{%
\RG{0.000}{-15.578}{10.9091}{rm}{49}%
}\hskip.1em\allowbreak{} 件\allowbreak{}\hskip.1em \RB{76.868}{5.317}{16.318}{%
\RG{0.000}{-0.005}{10.9091}{mi}{110}\RG{3.543}{-15.578}{10.9091}{mi}{110}\RG{6.549}{-1.641}{7.9701}{rm}{43}\RG{29.447}{-0.005}{10.9091}{mi}{110}\RG{33.536}{-15.578}{10.9091}{rm}{48}\RG{35.995}{-1.641}{7.9701}{sy}{0}\RG{65.923}{-0.005}{10.9091}{mi}{110}\RG{68.827}{-15.578}{10.9091}{rm}{48}\RG{72.471}{-1.641}{7.9701}{rm}{48}%
}\hskip.1em\allowbreak{}

\[
\RB{131.542}{8.679}{18.405}{%
\RG{0.000}{-0.003}{10.9091}{sy}{106}\RG{3.031}{-0.003}{10.9091}{mi}{97}\RG{8.797}{-0.003}{10.9091}{sy}{106}\RG{14.857}{-0.003}{10.9091}{rm}{61}\RG{26.355}{-0.003}{10.9091}{rm}{49}\RG{51.784}{-0.003}{10.9091}{rm}{49}\RG{81.228}{-0.003}{10.9091}{rm}{48}\RG{106.589}{-0.003}{10.9091}{mi}{110}\RG{115.561}{-0.003}{10.9091}{sy}{0}\RG{126.471}{-0.003}{10.9091}{rm}{49}\RG{0.000}{-15.178}{10.9091}{sy}{106}\RG{3.030}{-15.178}{10.9091}{mi}{97}\RG{8.797}{-15.178}{10.9091}{sy}{106}\RG{14.857}{-15.178}{10.9091}{mi}{62}\RG{26.373}{-15.178}{10.9091}{rm}{49}\RG{51.784}{-15.178}{10.9091}{rm}{49}\RG{71.291}{-15.178}{10.9091}{mi}{110}\RG{80.263}{-15.178}{10.9091}{sy}{0}\RG{91.172}{-15.178}{10.9091}{rm}{49}\RG{116.525}{-15.178}{10.9091}{rm}{48}%
}
\]
特别地，当\allowbreak{}\hskip.1em \RB{33.868}{7.769}{0.617}{%
\RG{0.000}{0.003}{10.9091}{mi}{97}\RG{8.797}{0.003}{10.9091}{rm}{61}\RG{20.312}{0.003}{10.9091}{sy}{6}\RG{28.797}{0.003}{10.9091}{rm}{49}%
}\hskip.1em\allowbreak{} 时，令

\[
\RB{115.593}{11.426}{3.229}{%
\RG{51.152}{4.502}{7.9701}{rm}{50}\RG{85.894}{4.502}{7.9701}{mi}{110}\RG{91.032}{4.502}{7.9701}{sy}{0}\RG{97.619}{4.502}{7.9701}{rm}{49}\RT{106.593}{4.502}{7.9701}{T}\RG{0.000}{-0.002}{10.9091}{mi}{118}\RG{8.710}{-0.002}{10.9091}{rm}{61}\RG{20.208}{-0.002}{10.9091}{rm}{40}\RG{24.441}{-0.002}{10.9091}{rm}{49}\RG{29.922}{-0.002}{10.9091}{mi}{59}\RG{34.755}{-0.002}{10.9091}{mi}{97}\RG{40.515}{-0.002}{10.9091}{mi}{59}\RG{45.359}{-0.002}{10.9091}{mi}{97}\RG{55.884}{-0.002}{10.9091}{mi}{59}\RG{60.717}{-0.002}{10.9091}{mi}{58}\RG{65.561}{-0.002}{10.9091}{mi}{58}\RG{70.405}{-0.002}{10.9091}{mi}{58}\RG{75.237}{-0.002}{10.9091}{mi}{59}\RG{80.081}{-0.002}{10.9091}{mi}{97}\RG{102.351}{-0.002}{10.9091}{rm}{41}\RG{112.878}{-0.002}{10.9091}{mi}{59}%
}
\]
则\allowbreak{}\hskip.1em \RB{33.995}{8.308}{0.623}{%
\RG{0.000}{-0.003}{10.9091}{mi}{65}\RG{11.212}{-0.003}{10.9091}{rm}{61}\RG{22.727}{-0.003}{10.9091}{mi}{118}\RG{28.389}{-0.003}{10.9091}{mi}{118}%
}\hskip.1em\allowbreak{} T，所以\allowbreak{}\hskip.1em \RB{8.366}{8.308}{0.503}{%
\RG{0.000}{-0.003}{10.9091}{mi}{65}%
}\hskip.1em\allowbreak{} 为秩一半正定阵，也直接核对了边界情形。\allowbreak{}\hskip.1em \RB{8.366}{7.992}{0.503}{%
\RG{0.000}{-0.003}{10.9091}{am}{3}%
}\hskip.1em\allowbreak{}

方法小结：矩阵元素沿行列呈等比变化时，可尝试相邻行、列的配对消元。求惯性指数只需合同化为对角阵，无需求出全部特征值。

题目来源：谢启鸿《复旦大学高等代数习题集（2023版）》，第117页。

% Original page 3.
\section*{2022S11}
设\allowbreak{}\hskip.1em \RB{43.654}{8.685}{3.760}{%
\RG{0.000}{0.003}{10.9091}{mi}{65}\RG{11.212}{0.003}{10.9091}{rm}{61}\RG{22.710}{0.003}{10.9091}{rm}{40}\RG{26.969}{0.003}{10.9091}{mi}{97}\RG{32.736}{-1.634}{7.9701}{mi}{105}\RG{35.613}{-1.634}{7.9701}{mi}{106}\RG{40.001}{0.003}{10.9091}{rm}{41}%
}\hskip.1em\allowbreak{} 为\allowbreak{}\hskip.1em \RB{6.729}{5.325}{0.617}{%
\RG{0.000}{0.003}{10.9091}{mi}{110}%
}\hskip.1em\allowbreak{} 阶正定实对称阵，\allowbreak{}\hskip.1em \RB{43.210}{8.685}{3.760}{%
\RG{0.000}{0.003}{10.9091}{mi}{66}\RG{11.852}{0.003}{10.9091}{rm}{61}\RG{23.350}{0.003}{10.9091}{rm}{40}\RG{27.610}{0.003}{10.9091}{mi}{98}\RG{32.291}{-1.634}{7.9701}{mi}{105}\RG{35.169}{-1.634}{7.9701}{mi}{106}\RG{39.557}{0.003}{10.9091}{rm}{41}%
}\hskip.1em\allowbreak{} 为\allowbreak{}\hskip.1em \RB{6.729}{5.325}{0.617}{%
\RG{0.000}{0.003}{10.9091}{mi}{110}%
}\hskip.1em\allowbreak{} 阶半正定实对称阵，且\allowbreak{}\hskip.1em \RB{31.010}{8.074}{2.214}{%
\RG{0.000}{0.003}{10.9091}{mi}{98}\RG{4.682}{-1.634}{7.9701}{mi}{105}\RG{7.559}{-1.634}{7.9701}{mi}{105}\RG{13.977}{0.003}{10.9091}{mi}{62}\RG{25.492}{0.003}{10.9091}{rm}{48}%
}\hskip.1em\allowbreak{} 。证明Hadamard乘积

\[
\RB{74.726}{8.685}{3.760}{%
\RG{0.000}{0.003}{10.9091}{mi}{65}\RG{10.606}{0.003}{10.9091}{sy}{14}\RG{18.485}{0.003}{10.9091}{mi}{66}\RG{30.337}{0.003}{10.9091}{rm}{61}\RG{41.836}{0.003}{10.9091}{rm}{40}\RG{46.095}{0.003}{10.9091}{mi}{97}\RG{59.127}{0.003}{10.9091}{mi}{98}\RG{71.074}{0.003}{10.9091}{rm}{41}\RG{51.861}{-1.634}{7.9701}{mi}{105}\RG{54.739}{-1.634}{7.9701}{mi}{106}\RG{63.809}{-1.634}{7.9701}{mi}{105}\RG{66.686}{-1.634}{7.9701}{mi}{106}%
}
\]
是正定实对称阵。

证明：由实对称矩阵的正交对角化定理，\allowbreak{}\hskip.1em \RB{8.747}{7.949}{0.502}{%
\RG{0.000}{-0.002}{10.9091}{mi}{66}%
}\hskip.1em\allowbreak{} 的特征值均非负，故可以取其实平方根，从而将\allowbreak{}\hskip.1em \RB{8.747}{7.949}{0.502}{%
\RG{0.000}{-0.002}{10.9091}{mi}{66}%
}\hskip.1em\allowbreak{} 写成

\[
\RB{102.489}{17.648}{13.937}{%
\RG{66.662}{13.633}{7.9701}{mi}{110}\RG{61.352}{10.360}{10.9092}{ex}{88}\RT{40.522}{4.500}{7.9701}{T}\RT{93.490}{4.500}{7.9701}{T}\RG{0.000}{-0.004}{10.9091}{mi}{66}\RG{11.852}{-0.004}{10.9091}{rm}{61}\RG{23.367}{-0.004}{10.9091}{mi}{67}\RG{31.931}{-0.004}{10.9091}{mi}{67}\RG{49.837}{-0.004}{10.9091}{rm}{61}\RG{78.928}{-0.004}{10.9091}{mi}{99}\RG{88.769}{-0.004}{10.9091}{mi}{99}\RG{99.775}{-0.004}{10.9091}{mi}{59}\RG{83.649}{-1.781}{7.9701}{mi}{107}\RG{93.490}{-2.780}{7.9701}{mi}{107}\RG{61.510}{-13.357}{7.9701}{mi}{107}\RG{66.132}{-13.357}{7.9701}{rm}{61}\RG{72.715}{-13.357}{7.9701}{rm}{49}%
}
\]
其中\allowbreak{}\hskip.1em \RB{9.469}{5.325}{2.354}{%
\RG{0.000}{0.003}{10.9091}{mi}{99}\RG{4.721}{-1.774}{7.9701}{mi}{107}%
}\hskip.1em\allowbreak{} 为\allowbreak{}\hskip.1em \RB{8.791}{8.194}{0.737}{%
\RG{0.000}{0.003}{10.9091}{mi}{67}%
}\hskip.1em\allowbreak{} 的第\allowbreak{}\hskip.1em \RB{6.042}{8.074}{0.617}{%
\RG{0.000}{0.003}{10.9091}{mi}{107}%
}\hskip.1em\allowbreak{} 列。记\allowbreak{}\hskip.1em \RB{43.742}{8.685}{3.224}{%
\RG{0.000}{0.003}{10.9091}{mi}{67}\RG{11.608}{0.003}{10.9091}{rm}{61}\RG{23.106}{0.003}{10.9091}{rm}{40}\RG{27.365}{0.003}{10.9091}{mi}{99}\RG{32.086}{-1.774}{7.9701}{mi}{105}\RG{34.963}{-1.774}{7.9701}{mi}{107}\RG{40.089}{0.003}{10.9091}{rm}{41}%
}\hskip.1em\allowbreak{} ，并令

\[
\RB{132.005}{9.292}{3.228}{%
\RG{0.000}{-0.001}{10.9091}{mi}{68}\RG{17.182}{-0.001}{10.9091}{rm}{61}\RT{28.697}{-0.001}{10.9091}{d}\RT{34.763}{-0.001}{10.9091}{i}\RT{37.795}{-0.001}{10.9091}{a}\RT{43.250}{-0.001}{10.9091}{g}\RG{48.705}{-0.001}{10.9091}{rm}{40}\RG{52.947}{-0.001}{10.9091}{mi}{99}\RG{67.023}{-0.001}{10.9091}{mi}{59}\RG{71.855}{-0.001}{10.9091}{mi}{99}\RG{85.946}{-0.001}{10.9091}{mi}{59}\RG{90.778}{-0.001}{10.9091}{mi}{58}\RG{95.622}{-0.001}{10.9091}{mi}{58}\RG{100.466}{-0.001}{10.9091}{mi}{58}\RG{105.299}{-0.001}{10.9091}{mi}{59}\RG{110.142}{-0.001}{10.9091}{mi}{99}\RG{125.168}{-0.001}{10.9091}{rm}{41}\RG{129.410}{-0.001}{10.9091}{mi}{58}\RG{9.032}{-1.778}{7.9701}{mi}{107}\RG{57.668}{-1.778}{7.9701}{rm}{49}\RG{61.902}{-1.778}{7.9701}{mi}{107}\RG{76.592}{-1.778}{7.9701}{rm}{50}\RG{80.827}{-1.778}{7.9701}{mi}{107}\RG{114.910}{-1.778}{7.9701}{mi}{110}\RG{120.043}{-1.778}{7.9701}{mi}{107}%
}
\]
比较第\allowbreak{}\hskip.1em \RB{21.619}{8.687}{3.222}{%
\RG{0.000}{0.005}{10.9091}{rm}{40}\RG{4.242}{0.005}{10.9091}{mi}{105}\RG{7.995}{0.005}{10.9091}{mi}{59}\RG{12.828}{0.005}{10.9091}{mi}{106}\RG{17.966}{0.005}{10.9091}{rm}{41}%
}\hskip.1em\allowbreak{} 个元素，有

\[
\RB{194.139}{19.593}{17.067}{%
\RG{0.000}{18.656}{10.9092}{ex}{32}\RG{62.699}{18.656}{10.9092}{ex}{33}\RG{13.947}{13.638}{7.9701}{mi}{110}\RG{113.305}{13.638}{7.9701}{mi}{110}\RG{8.637}{10.365}{10.9092}{ex}{88}\RG{107.995}{10.365}{10.9092}{ex}{88}\RG{26.213}{0.002}{10.9091}{mi}{68}\RG{40.364}{0.002}{10.9091}{mi}{65}\RG{48.546}{0.002}{10.9091}{mi}{68}\RG{81.630}{0.002}{10.9091}{rm}{61}\RG{93.146}{0.002}{10.9091}{mi}{97}\RG{125.571}{0.002}{10.9091}{mi}{99}\RG{138.295}{0.002}{10.9091}{mi}{99}\RG{155.050}{0.002}{10.9091}{rm}{61}\RG{166.565}{0.002}{10.9091}{mi}{97}\RG{179.598}{0.002}{10.9091}{mi}{98}\RG{191.545}{0.002}{10.9091}{mi}{58}\RG{98.912}{-1.635}{7.9701}{mi}{105}\RG{101.789}{-1.635}{7.9701}{mi}{106}\RG{172.332}{-1.635}{7.9701}{mi}{105}\RG{175.210}{-1.635}{7.9701}{mi}{106}\RG{184.280}{-1.635}{7.9701}{mi}{105}\RG{187.157}{-1.635}{7.9701}{mi}{106}\RG{35.244}{-1.776}{7.9701}{mi}{107}\RG{57.578}{-1.776}{7.9701}{mi}{107}\RG{130.292}{-1.776}{7.9701}{mi}{105}\RG{133.169}{-1.776}{7.9701}{mi}{107}\RG{143.016}{-1.776}{7.9701}{mi}{106}\RG{146.890}{-1.776}{7.9701}{mi}{107}\RG{8.794}{-13.351}{7.9701}{mi}{107}\RG{13.416}{-13.352}{7.9701}{rm}{61}\RG{19.999}{-13.352}{7.9701}{rm}{49}\RG{108.153}{-13.351}{7.9701}{mi}{107}\RG{112.775}{-13.352}{7.9701}{rm}{61}\RG{119.358}{-13.352}{7.9701}{rm}{49}\RG{71.335}{-14.941}{7.9701}{mi}{105}\RG{74.212}{-14.941}{7.9701}{mi}{106}%
}
\]
因此

\[
\RB{98.508}{17.656}{13.928}{%
\RG{47.162}{13.641}{7.9701}{mi}{110}\RG{41.853}{10.369}{10.9092}{ex}{88}\RG{0.000}{0.005}{10.9091}{mi}{65}\RG{10.606}{0.005}{10.9091}{sy}{14}\RG{18.485}{0.005}{10.9091}{mi}{66}\RG{30.337}{0.005}{10.9091}{rm}{61}\RG{59.428}{0.005}{10.9091}{mi}{68}\RG{73.581}{0.005}{10.9091}{mi}{65}\RG{81.762}{0.005}{10.9091}{mi}{68}\RG{95.914}{0.005}{10.9091}{mi}{58}\RG{68.461}{-1.772}{7.9701}{mi}{107}\RG{90.794}{-1.772}{7.9701}{mi}{107}\RG{42.010}{-13.348}{7.9701}{mi}{107}\RG{46.632}{-13.348}{7.9701}{rm}{61}\RG{53.216}{-13.348}{7.9701}{rm}{49}%
}
\]
任取\allowbreak{}\hskip.1em \RB{59.454}{8.679}{3.230}{%
\RG{0.000}{-0.003}{10.9091}{mi}{120}\RG{9.265}{-0.003}{10.9091}{sy}{50}\RG{19.568}{-0.003}{10.9091}{bb}{82}\RG{27.447}{3.956}{7.9701}{mi}{110}\RG{35.507}{-0.003}{10.9091}{sy}{110}\RG{43.384}{-0.003}{10.9091}{sy}{102}\RG{48.841}{-0.003}{10.9091}{rm}{48}\RG{54.295}{-0.003}{10.9091}{sy}{103}%
}\hskip.1em\allowbreak{} 。由于\allowbreak{}\hskip.1em \RB{8.366}{8.308}{0.503}{%
\RG{0.000}{-0.003}{10.9091}{mi}{65}%
}\hskip.1em\allowbreak{} 正定，

\[
\RB{179.654}{17.647}{13.937}{%
\RG{74.402}{13.632}{7.9701}{mi}{110}\RG{69.092}{10.360}{10.9092}{ex}{88}\RT{6.235}{4.500}{7.9701}{T}\RT{113.721}{4.500}{7.9701}{T}\RG{0.000}{-0.004}{10.9091}{mi}{120}\RG{12.519}{-0.004}{10.9091}{rm}{40}\RG{16.762}{-0.004}{10.9091}{mi}{65}\RG{27.368}{-0.004}{10.9091}{sy}{14}\RG{35.247}{-0.004}{10.9091}{mi}{66}\RG{44.069}{-0.004}{10.9091}{rm}{41}\RG{48.311}{-0.004}{10.9091}{mi}{120}\RG{57.576}{-0.004}{10.9091}{rm}{61}\RG{84.850}{-0.004}{10.9091}{rm}{40}\RG{89.092}{-0.004}{10.9091}{mi}{68}\RG{103.244}{-0.004}{10.9091}{mi}{120}\RG{109.479}{-0.004}{10.9091}{rm}{41}\RG{120.006}{-0.004}{10.9091}{mi}{65}\RG{128.187}{-0.004}{10.9091}{rm}{40}\RG{132.430}{-0.004}{10.9091}{mi}{68}\RG{146.581}{-0.004}{10.9091}{mi}{120}\RG{152.816}{-0.004}{10.9091}{rm}{41}\RG{160.089}{-0.004}{10.9091}{sy}{21}\RG{171.604}{-0.004}{10.9091}{rm}{48}\RG{177.059}{-0.004}{10.9091}{mi}{58}\RG{98.124}{-1.781}{7.9701}{mi}{107}\RG{141.462}{-1.781}{7.9701}{mi}{107}\RG{69.250}{-13.357}{7.9701}{mi}{107}\RG{73.871}{-13.357}{7.9701}{rm}{61}\RG{80.454}{-13.357}{7.9701}{rm}{49}%
}
\]
要使不等式严格，只需证明至少有一个\allowbreak{}\hskip.1em \RB{20.401}{7.951}{2.358}{%
\RG{0.000}{-0.000}{10.9091}{mi}{68}\RG{9.032}{-1.778}{7.9701}{mi}{107}\RG{14.152}{-0.000}{10.9091}{mi}{120}%
}\hskip.1em\allowbreak{} 不为零。取\allowbreak{}\hskip.1em \RB{29.664}{8.311}{2.857}{%
\RG{0.000}{-0.000}{10.9091}{mi}{120}\RG{6.235}{-1.637}{7.9701}{mi}{105}\RG{12.647}{-0.000}{10.9091}{sy}{54}\RG{12.647}{-0.000}{10.9091}{rm}{61}\RG{24.145}{-0.000}{10.9091}{rm}{48}%
}\hskip.1em\allowbreak{} 的一个指标\allowbreak{}\hskip.1em \RB{3.696}{7.711}{0.620}{%
\RG{0.000}{-0.000}{10.9091}{mi}{105}%
}\hskip.1em\allowbreak{} 。由

\[
\RB{78.506}{17.656}{13.929}{%
\RG{50.801}{13.641}{7.9701}{mi}{110}\RG{45.492}{10.368}{10.9092}{ex}{88}\RG{67.789}{5.312}{7.9701}{mi}{59}\RG{70.142}{5.312}{7.9701}{rm}{50}\RG{0.000}{0.004}{10.9091}{rm}{48}\RG{8.485}{0.004}{10.9091}{mi}{60}\RG{19.983}{0.004}{10.9091}{mi}{98}\RG{33.977}{0.004}{10.9091}{rm}{61}\RG{63.068}{0.004}{10.9091}{mi}{99}\RG{75.792}{0.004}{10.9091}{mi}{59}\RG{24.682}{-1.632}{7.9701}{mi}{105}\RG{27.559}{-1.632}{7.9701}{mi}{105}\RG{67.789}{-3.518}{7.9701}{mi}{105}\RG{70.666}{-3.518}{7.9701}{mi}{107}\RG{45.649}{-13.349}{7.9701}{mi}{107}\RG{50.271}{-13.349}{7.9701}{rm}{61}\RG{56.855}{-13.349}{7.9701}{rm}{49}%
}
\]
可选到\allowbreak{}\hskip.1em \RB{10.076}{8.069}{2.306}{%
\RG{0.000}{-0.002}{10.9091}{mi}{107}\RG{5.679}{-1.638}{7.9701}{rm}{48}%
}\hskip.1em\allowbreak{} 使\allowbreak{}\hskip.1em \RB{36.710}{8.309}{3.512}{%
\RG{0.000}{-0.002}{10.9091}{mi}{99}\RG{4.721}{-1.779}{7.9701}{mi}{105}\RG{7.598}{-1.779}{7.9701}{mi}{107}\RG{12.014}{-2.886}{5.9776}{rm}{48}\RG{19.693}{-0.002}{10.9091}{sy}{54}\RG{19.693}{-0.002}{10.9091}{rm}{61}\RG{31.191}{-0.002}{10.9091}{rm}{48}%
}\hskip.1em\allowbreak{} 。于是

\[
\RB{99.731}{8.685}{3.507}{%
\RG{0.000}{0.003}{10.9091}{rm}{40}\RG{4.242}{0.003}{10.9091}{mi}{68}\RG{22.333}{0.003}{10.9091}{mi}{120}\RG{28.568}{0.003}{10.9091}{rm}{41}\RG{39.222}{0.003}{10.9091}{rm}{61}\RG{50.737}{0.003}{10.9091}{mi}{99}\RG{67.400}{0.003}{10.9091}{mi}{120}\RG{80.047}{0.003}{10.9091}{sy}{54}\RG{80.047}{0.003}{10.9091}{rm}{61}\RG{91.545}{0.003}{10.9091}{rm}{48}\RG{97.017}{0.003}{10.9091}{mi}{59}\RG{32.810}{-1.633}{7.9701}{mi}{105}\RG{73.635}{-1.633}{7.9701}{mi}{105}\RG{13.274}{-1.774}{7.9701}{mi}{107}\RG{55.458}{-1.774}{7.9701}{mi}{105}\RG{58.336}{-1.774}{7.9701}{mi}{107}\RG{17.684}{-2.881}{5.9776}{rm}{48}\RG{62.751}{-2.881}{5.9776}{rm}{48}%
}
\]
从而\allowbreak{}\hskip.1em \RB{32.221}{8.684}{3.508}{%
\RG{0.000}{0.002}{10.9091}{rm}{40}\RG{4.242}{0.002}{10.9091}{mi}{68}\RG{13.274}{-1.775}{7.9701}{mi}{107}\RG{17.684}{-2.882}{5.9776}{rm}{48}\RG{22.333}{0.002}{10.9091}{mi}{120}\RG{28.568}{0.002}{10.9091}{rm}{41}%
}\hskip.1em\allowbreak{} T\allowbreak{}\hskip.1em \RB{61.057}{8.684}{3.508}{%
\RG{0.000}{0.002}{10.9091}{mi}{65}\RG{8.182}{0.002}{10.9091}{rm}{40}\RG{12.425}{0.002}{10.9091}{mi}{68}\RG{21.457}{-1.775}{7.9701}{mi}{107}\RG{25.866}{-2.882}{5.9776}{rm}{48}\RG{30.515}{0.002}{10.9091}{mi}{120}\RG{36.750}{0.002}{10.9091}{rm}{41}\RG{44.023}{0.002}{10.9091}{mi}{62}\RG{55.538}{0.002}{10.9091}{rm}{48}%
}\hskip.1em\allowbreak{} ，其余各项均非负。因此

\[
\RB{77.141}{11.427}{3.229}{%
\RT{6.235}{4.503}{7.9701}{T}\RG{0.000}{-0.002}{10.9091}{mi}{120}\RG{12.519}{-0.002}{10.9091}{rm}{40}\RG{16.762}{-0.002}{10.9091}{mi}{65}\RG{27.368}{-0.002}{10.9091}{sy}{14}\RG{35.246}{-0.002}{10.9091}{mi}{66}\RG{44.068}{-0.002}{10.9091}{rm}{41}\RG{48.311}{-0.002}{10.9091}{mi}{120}\RG{57.562}{-0.002}{10.9091}{mi}{62}\RG{69.092}{-0.002}{10.9091}{rm}{48}\RG{74.547}{-0.002}{10.9091}{mi}{58}%
}
\]
又因\allowbreak{}\hskip.1em \RB{21.762}{8.315}{2.602}{%
\RG{0.000}{0.004}{10.9091}{mi}{65}\RG{8.182}{0.004}{10.9091}{mi}{59}\RG{13.015}{0.004}{10.9091}{mi}{66}%
}\hskip.1em\allowbreak{} 均对称，\allowbreak{}\hskip.1em \RB{27.233}{8.315}{0.500}{%
\RG{0.000}{0.004}{10.9091}{mi}{65}\RG{10.606}{0.004}{10.9091}{sy}{14}\RG{18.485}{0.004}{10.9091}{mi}{66}%
}\hskip.1em\allowbreak{} 也对称，故结论成立。\allowbreak{}\hskip.1em \RB{8.366}{7.999}{0.500}{%
\RG{0.000}{0.004}{10.9091}{am}{3}%
}\hskip.1em\allowbreak{}

条件\allowbreak{}\hskip.1em \RB{31.010}{8.072}{2.216}{%
\RG{0.000}{0.001}{10.9091}{mi}{98}\RG{4.682}{-1.636}{7.9701}{mi}{105}\RG{7.559}{-1.636}{7.9701}{mi}{105}\RG{13.976}{0.001}{10.9091}{mi}{62}\RG{25.492}{0.001}{10.9091}{rm}{48}%
}\hskip.1em\allowbreak{} 不能省略：若某个\allowbreak{}\hskip.1em \RB{30.994}{8.072}{2.216}{%
\RG{0.000}{0.001}{10.9091}{mi}{98}\RG{4.682}{-1.636}{7.9701}{mi}{105}\RG{7.559}{-1.636}{7.9701}{mi}{105}\RG{13.977}{0.001}{10.9091}{rm}{61}\RG{25.475}{0.001}{10.9091}{rm}{48}%
}\hskip.1em\allowbreak{} ，则

\[
\RB{115.855}{11.425}{3.280}{%
\RT{5.080}{4.501}{7.9701}{T}\RG{0.000}{-0.003}{10.9091}{mi}{101}\RG{11.364}{-0.003}{10.9091}{rm}{40}\RG{15.607}{-0.003}{10.9091}{mi}{65}\RG{26.213}{-0.003}{10.9091}{sy}{14}\RG{34.092}{-0.003}{10.9091}{mi}{66}\RG{42.914}{-0.003}{10.9091}{rm}{41}\RG{47.156}{-0.003}{10.9091}{mi}{101}\RG{58.647}{-0.003}{10.9091}{rm}{61}\RG{70.163}{-0.003}{10.9091}{mi}{97}\RG{82.194}{-0.003}{10.9091}{mi}{98}\RG{96.171}{-0.003}{10.9091}{rm}{61}\RG{107.669}{-0.003}{10.9091}{rm}{48}\RG{113.140}{-0.003}{10.9091}{mi}{59}\RG{52.236}{-1.640}{7.9701}{mi}{105}\RG{75.930}{-1.640}{7.9701}{mi}{105}\RG{78.807}{-1.640}{7.9701}{mi}{105}\RG{86.876}{-1.640}{7.9701}{mi}{105}\RG{89.753}{-1.640}{7.9701}{mi}{105}\RG{5.080}{-2.700}{7.9701}{mi}{105}%
}
\]
因而\allowbreak{}\hskip.1em \RB{27.233}{8.314}{0.500}{%
\RG{0.000}{0.003}{10.9091}{mi}{65}\RG{10.606}{0.003}{10.9091}{sy}{14}\RG{18.485}{0.003}{10.9091}{mi}{66}%
}\hskip.1em\allowbreak{} 不可能正定。

方法小结：将半正定阵分解为秩一矩阵之和，把Hadamard乘积改写为对角合同变换之和；对角元严格为正保证其中至少一项严格为正。

题目来源：谢启鸿《复旦大学高等代数习题集（2023版）》，第117页。

% Original page 4.
\section*{2022S12}
设\allowbreak{}\hskip.1em \RB{8.366}{8.314}{0.500}{%
\RG{0.000}{0.003}{10.9091}{mi}{65}%
}\hskip.1em\allowbreak{} 为\allowbreak{}\hskip.1em \RB{6.729}{5.325}{0.617}{%
\RG{0.000}{0.003}{10.9091}{mi}{110}%
}\hskip.1em\allowbreak{} 阶实对称阵，\allowbreak{}\hskip.1em \RB{8.747}{7.954}{0.500}{%
\RG{0.000}{0.003}{10.9091}{mi}{66}%
}\hskip.1em\allowbreak{} 为\allowbreak{}\hskip.1em \RB{6.729}{5.325}{0.617}{%
\RG{0.000}{0.003}{10.9091}{mi}{110}%
}\hskip.1em\allowbreak{} 阶半正定实对称阵，满足

\[
\RB{79.602}{9.292}{3.228}{%
\RT{0.000}{-0.001}{10.9091}{d}\RT{6.066}{-0.001}{10.9091}{e}\RT{10.909}{-0.001}{10.9091}{t}\RG{15.153}{-0.001}{10.9091}{rm}{40}\RG{19.395}{-0.001}{10.9091}{mi}{65}\RG{30.001}{-0.001}{10.9091}{rm}{43}\RT{40.910}{-0.001}{10.9091}{i}\RG{43.943}{-0.001}{10.9091}{mi}{66}\RG{52.765}{-0.001}{10.9091}{rm}{41}\RG{60.020}{-0.001}{10.9091}{rm}{61}\RG{71.529}{-0.001}{10.9091}{rm}{48}\RG{77.008}{-0.001}{10.9091}{mi}{58}%
}
\]
证明：存在非零实列向量\allowbreak{}\hskip.1em \RB{33.609}{7.979}{0.931}{%
\RG{0.000}{0.005}{10.9091}{mi}{11}\RG{10.050}{0.005}{10.9091}{sy}{50}\RG{20.353}{0.005}{10.9091}{bb}{82}\RG{28.231}{3.964}{7.9701}{mi}{110}%
}\hskip.1em\allowbreak{} ，使得\allowbreak{}\hskip.1em \RB{65.634}{8.316}{0.735}{%
\RG{0.000}{0.005}{10.9091}{mi}{65}\RG{8.182}{0.005}{10.9091}{mi}{11}\RG{18.231}{0.005}{10.9091}{rm}{61}\RG{29.746}{0.005}{10.9091}{mi}{66}\RG{38.561}{0.005}{10.9091}{mi}{11}\RG{48.617}{0.005}{10.9091}{rm}{61}\RG{60.116}{0.005}{10.9091}{rm}{48}%
}\hskip.1em\allowbreak{} 。证明：由行列式为零，存在\allowbreak{}\hskip.1em \RB{58.772}{8.683}{3.226}{%
\RG{0.000}{0.001}{10.9091}{mi}{122}\RG{8.583}{0.001}{10.9091}{sy}{50}\RG{18.886}{0.001}{10.9091}{bb}{67}\RG{26.765}{3.960}{7.9701}{mi}{110}\RG{34.825}{0.001}{10.9091}{sy}{110}\RG{42.702}{0.001}{10.9091}{sy}{102}\RG{48.159}{0.001}{10.9091}{rm}{48}\RG{53.613}{0.001}{10.9091}{sy}{103}%
}\hskip.1em\allowbreak{} ，使

\[
\RB{70.003}{9.290}{3.230}{%
\RG{0.000}{-0.003}{10.9091}{rm}{40}\RG{4.242}{-0.003}{10.9091}{mi}{65}\RG{14.848}{-0.003}{10.9091}{rm}{43}\RT{25.757}{-0.003}{10.9091}{i}\RG{28.790}{-0.003}{10.9091}{mi}{66}\RG{37.612}{-0.003}{10.9091}{rm}{41}\RG{41.855}{-0.003}{10.9091}{mi}{122}\RG{50.438}{-0.003}{10.9091}{rm}{61}\RG{61.937}{-0.003}{10.9091}{rm}{48}\RG{67.409}{-0.003}{10.9091}{mi}{58}%
}
\]
左乘\allowbreak{}\hskip.1em \RB{9.719}{8.168}{0.617}{%
\RG{0.000}{0.003}{10.9091}{mi}{122}\RG{5.553}{3.962}{7.9701}{sy}{3}%
}\hskip.1em\allowbreak{} ，其中\allowbreak{}\hskip.1em \RB{9.719}{8.168}{0.617}{%
\RG{0.000}{0.003}{10.9091}{mi}{122}\RG{5.553}{3.962}{7.9701}{sy}{3}%
}\hskip.1em\allowbreak{} 表示共轭转置，得

\[
\RB{89.460}{9.292}{2.617}{%
\RG{5.553}{4.504}{7.9701}{sy}{3}\RG{47.758}{4.504}{7.9701}{sy}{3}\RG{0.000}{-0.001}{10.9091}{mi}{122}\RG{10.286}{-0.001}{10.9091}{mi}{65}\RG{18.468}{-0.001}{10.9091}{mi}{122}\RG{26.445}{-0.001}{10.9091}{rm}{43}\RT{37.354}{-0.001}{10.9091}{i}\RG{42.205}{-0.001}{10.9091}{mi}{122}\RG{52.490}{-0.001}{10.9091}{mi}{66}\RG{61.305}{-0.001}{10.9091}{mi}{122}\RG{69.896}{-0.001}{10.9091}{rm}{61}\RG{81.394}{-0.001}{10.9091}{rm}{48}\RG{86.866}{-0.001}{10.9091}{mi}{58}%
}
\]
因为实对称阵也是Hermite阵，\allowbreak{}\hskip.1em \RB{24.063}{8.316}{0.615}{%
\RG{0.000}{0.005}{10.9091}{mi}{122}\RG{5.553}{3.964}{7.9701}{sy}{3}\RG{10.286}{0.005}{10.9091}{mi}{65}\RG{18.468}{0.005}{10.9091}{mi}{122}%
}\hskip.1em\allowbreak{} 和\allowbreak{}\hskip.1em \RB{24.695}{8.170}{0.615}{%
\RG{0.000}{0.005}{10.9091}{mi}{122}\RG{5.553}{3.964}{7.9701}{sy}{3}\RG{10.286}{0.005}{10.9091}{mi}{66}\RG{19.101}{0.005}{10.9091}{mi}{122}%
}\hskip.1em\allowbreak{} 均为实数。比较虚部，得到

\[
\RB{47.255}{8.712}{0.739}{%
\RG{5.553}{4.506}{7.9701}{sy}{3}\RG{0.000}{0.001}{10.9091}{mi}{122}\RG{10.285}{0.001}{10.9091}{mi}{66}\RG{19.100}{0.001}{10.9091}{mi}{122}\RG{27.690}{0.001}{10.9091}{rm}{61}\RG{39.188}{0.001}{10.9091}{rm}{48}\RG{44.660}{0.001}{10.9091}{mi}{58}%
}
\]
下面利用\allowbreak{}\hskip.1em \RB{8.747}{7.950}{0.501}{%
\RG{0.000}{-0.001}{10.9091}{mi}{66}%
}\hskip.1em\allowbreak{} 的半正定性，将二次型等于零转化为向量属于核。由正交对角化，可写

\[
\RB{249.455}{11.433}{3.758}{%
\RT{118.685}{4.509}{7.9701}{T}\RT{160.260}{4.509}{7.9701}{T}\RG{0.000}{0.005}{10.9091}{mi}{66}\RG{11.852}{0.005}{10.9091}{rm}{61}\RG{23.367}{0.005}{10.9091}{mi}{81}\RT{33.810}{0.005}{10.9091}{d}\RT{39.876}{0.005}{10.9091}{i}\RT{42.909}{0.005}{10.9091}{a}\RT{48.363}{0.005}{10.9091}{g}\RG{53.817}{0.005}{10.9091}{rm}{40}\RG{58.059}{0.005}{10.9091}{mi}{22}\RG{69.366}{0.005}{10.9091}{mi}{59}\RG{74.198}{0.005}{10.9091}{mi}{58}\RG{79.042}{0.005}{10.9091}{mi}{58}\RG{83.875}{0.005}{10.9091}{mi}{58}\RG{88.718}{0.005}{10.9091}{mi}{59}\RG{93.562}{0.005}{10.9091}{mi}{22}\RG{105.818}{0.005}{10.9091}{rm}{41}\RG{110.060}{0.005}{10.9091}{mi}{81}\RG{124.969}{0.005}{10.9091}{mi}{59}\RG{151.620}{0.005}{10.9091}{mi}{81}\RG{166.544}{0.005}{10.9091}{mi}{81}\RG{178.199}{0.005}{10.9091}{rm}{61}\RG{189.715}{0.005}{10.9091}{mi}{73}\RG{200.147}{0.005}{10.9091}{mi}{59}\RG{215.889}{0.005}{10.9091}{mi}{22}\RG{229.890}{0.005}{10.9091}{sy}{21}\RG{241.405}{0.005}{10.9091}{rm}{48}\RG{246.860}{0.005}{10.9091}{mi}{58}\RG{64.634}{-1.632}{7.9701}{rm}{49}\RG{100.182}{-1.632}{7.9701}{mi}{110}\RG{194.510}{-1.632}{7.9701}{mi}{110}\RG{222.477}{-1.632}{7.9701}{mi}{106}%
}
\]
令\allowbreak{}\hskip.1em \RB{31.222}{8.192}{2.615}{%
\RG{0.000}{0.001}{10.9091}{mi}{119}\RG{11.134}{0.001}{10.9091}{rm}{61}\RG{22.649}{0.001}{10.9091}{mi}{81}%
}\hskip.1em\allowbreak{} T\allowbreak{}\hskip.1em \RB{5.595}{5.323}{0.619}{%
\RG{0.000}{0.001}{10.9091}{mi}{122}%
}\hskip.1em\allowbreak{} ，则

\[
\RB{113.317}{17.656}{15.209}{%
\RG{64.516}{13.641}{7.9701}{mi}{110}\RG{59.205}{10.368}{10.9092}{ex}{88}\RG{25.553}{4.509}{7.9701}{sy}{3}\RG{105.990}{4.509}{7.9701}{rm}{50}\RG{0.000}{0.005}{10.9091}{rm}{48}\RG{8.476}{0.005}{10.9091}{rm}{61}\RG{20.000}{0.005}{10.9091}{mi}{122}\RG{30.285}{0.005}{10.9091}{mi}{66}\RG{39.100}{0.005}{10.9091}{mi}{122}\RG{47.690}{0.005}{10.9091}{rm}{61}\RG{76.782}{0.005}{10.9091}{mi}{22}\RG{87.738}{0.005}{10.9091}{sy}{106}\RG{90.768}{0.005}{10.9091}{mi}{119}\RG{102.960}{0.005}{10.9091}{sy}{106}\RG{110.723}{0.005}{10.9091}{mi}{58}\RG{83.355}{-1.632}{7.9701}{mi}{106}\RG{98.578}{-1.632}{7.9701}{mi}{106}\RG{59.732}{-13.083}{7.9701}{mi}{106}\RG{63.617}{-13.083}{7.9701}{rm}{61}\RG{70.200}{-13.083}{7.9701}{rm}{49}%
}
\]
每个加项均非负，故\allowbreak{}\hskip.1em \RB{53.988}{9.758}{3.763}{%
\RG{0.000}{-0.000}{10.9091}{mi}{22}\RG{6.574}{-1.637}{7.9701}{mi}{106}\RG{10.956}{-0.000}{10.9091}{sy}{106}\RG{13.986}{-0.000}{10.9091}{mi}{119}\RG{21.796}{-1.637}{7.9701}{mi}{106}\RG{26.178}{-0.000}{10.9091}{sy}{106}\RG{29.209}{3.958}{7.9701}{rm}{50}\RG{36.971}{-0.000}{10.9091}{rm}{61}\RG{48.470}{-0.000}{10.9091}{rm}{48}%
}\hskip.1em\allowbreak{} 。若\allowbreak{}\hskip.1em \RB{31.019}{7.766}{3.763}{%
\RG{0.000}{-0.000}{10.9091}{mi}{22}\RG{6.573}{-1.637}{7.9701}{mi}{106}\RG{13.985}{-0.000}{10.9091}{mi}{62}\RG{25.501}{-0.000}{10.9091}{rm}{48}%
}\hskip.1em\allowbreak{} ，便有\allowbreak{}\hskip.1em \RB{32.239}{7.766}{3.763}{%
\RG{0.000}{-0.000}{10.9091}{mi}{119}\RG{7.810}{-1.637}{7.9701}{mi}{106}\RG{15.222}{-0.000}{10.9091}{rm}{61}\RG{26.720}{-0.000}{10.9091}{rm}{48}%
}\hskip.1em\allowbreak{} ；若\allowbreak{}\hskip.1em \RB{31.003}{7.766}{3.763}{%
\RG{0.000}{-0.000}{10.9091}{mi}{22}\RG{6.573}{-1.637}{7.9701}{mi}{106}\RG{13.986}{-0.000}{10.9091}{rm}{61}\RG{25.484}{-0.000}{10.9091}{rm}{48}%
}\hskip.1em\allowbreak{} ，也有\allowbreak{}\hskip.1em \RB{43.195}{7.766}{3.763}{%
\RG{0.000}{-0.000}{10.9091}{mi}{22}\RG{6.574}{-1.637}{7.9701}{mi}{106}\RG{10.956}{-0.000}{10.9091}{mi}{119}\RG{18.766}{-1.637}{7.9701}{mi}{106}\RG{26.178}{-0.000}{10.9091}{rm}{61}\RG{37.676}{-0.000}{10.9091}{rm}{48}%
}\hskip.1em\allowbreak{} 。因此

\[
\RB{167.991}{9.289}{3.231}{%
\RT{0.000}{-0.004}{10.9091}{d}\RT{6.066}{-0.004}{10.9091}{i}\RT{9.098}{-0.004}{10.9091}{a}\RT{14.553}{-0.004}{10.9091}{g}\RG{20.008}{-0.004}{10.9091}{rm}{40}\RG{24.250}{-0.004}{10.9091}{mi}{22}\RG{35.556}{-0.004}{10.9091}{mi}{59}\RG{40.389}{-0.004}{10.9091}{mi}{58}\RG{45.233}{-0.004}{10.9091}{mi}{58}\RG{50.066}{-0.004}{10.9091}{mi}{58}\RG{54.909}{-0.004}{10.9091}{mi}{59}\RG{59.753}{-0.004}{10.9091}{mi}{22}\RG{72.009}{-0.004}{10.9091}{rm}{41}\RG{76.251}{-0.004}{10.9091}{mi}{119}\RG{87.385}{-0.004}{10.9091}{rm}{61}\RG{98.883}{-0.004}{10.9091}{rm}{48}\RG{104.354}{-0.004}{10.9091}{mi}{59}\RG{131.005}{-0.004}{10.9091}{mi}{66}\RG{139.820}{-0.004}{10.9091}{mi}{122}\RG{148.426}{-0.004}{10.9091}{rm}{61}\RG{159.925}{-0.004}{10.9091}{rm}{48}\RG{165.396}{-0.004}{10.9091}{mi}{58}\RG{30.823}{-1.641}{7.9701}{rm}{49}\RG{66.372}{-1.641}{7.9701}{mi}{110}%
}
\]
再代回\allowbreak{}\hskip.1em \RB{23.214}{8.684}{3.225}{%
\RG{0.000}{0.002}{10.9091}{rm}{40}\RG{4.242}{0.002}{10.9091}{mi}{65}\RG{14.848}{0.002}{10.9091}{rm}{43}%
}\hskip.1em\allowbreak{} i\allowbreak{}\hskip.1em \RB{38.665}{8.684}{3.225}{%
\RG{0.000}{0.002}{10.9091}{mi}{66}\RG{8.822}{0.002}{10.9091}{rm}{41}\RG{13.065}{0.002}{10.9091}{mi}{122}\RG{21.648}{0.002}{10.9091}{rm}{61}\RG{33.146}{0.002}{10.9091}{rm}{48}%
}\hskip.1em\allowbreak{} ，可得\allowbreak{}\hskip.1em \RB{33.782}{8.313}{0.738}{%
\RG{0.000}{0.002}{10.9091}{mi}{65}\RG{8.182}{0.002}{10.9091}{mi}{122}\RG{16.765}{0.002}{10.9091}{rm}{61}\RG{28.263}{0.002}{10.9091}{rm}{48}%
}\hskip.1em\allowbreak{} 。

最后写\allowbreak{}\hskip.1em \RB{37.003}{6.858}{1.407}{%
\RG{0.000}{-0.002}{10.9091}{mi}{122}\RG{8.583}{-0.002}{10.9091}{rm}{61}\RG{20.098}{-0.002}{10.9091}{mi}{117}\RG{28.637}{-0.002}{10.9091}{rm}{43}%
}\hskip.1em\allowbreak{} i\allowbreak{}\hskip.1em \RB{5.606}{5.320}{0.622}{%
\RG{0.000}{-0.002}{10.9091}{mi}{118}%
}\hskip.1em\allowbreak{} ，其中\allowbreak{}\hskip.1em \RB{43.363}{7.972}{2.607}{%
\RG{0.000}{-0.002}{10.9091}{mi}{117}\RG{6.240}{-0.002}{10.9091}{mi}{59}\RG{11.073}{-0.002}{10.9091}{mi}{118}\RG{19.803}{-0.002}{10.9091}{sy}{50}\RG{30.106}{-0.002}{10.9091}{bb}{82}\RG{37.985}{3.957}{7.9701}{mi}{110}%
}\hskip.1em\allowbreak{} 。由于\allowbreak{}\hskip.1em \RB{21.762}{8.309}{2.607}{%
\RG{0.000}{-0.002}{10.9091}{mi}{65}\RG{8.182}{-0.002}{10.9091}{mi}{59}\RG{13.015}{-0.002}{10.9091}{mi}{66}%
}\hskip.1em\allowbreak{} 的元素均为实数，分别比较实部与虚部，得到

\[
\RB{124.088}{8.315}{0.736}{%
\RG{0.000}{0.004}{10.9091}{mi}{65}\RG{8.182}{0.004}{10.9091}{mi}{117}\RG{17.457}{0.004}{10.9091}{rm}{61}\RG{28.972}{0.004}{10.9091}{mi}{65}\RG{37.154}{0.004}{10.9091}{mi}{118}\RG{45.863}{0.004}{10.9091}{rm}{61}\RG{57.378}{0.004}{10.9091}{mi}{66}\RG{66.193}{0.004}{10.9091}{mi}{117}\RG{75.477}{0.004}{10.9091}{rm}{61}\RG{86.992}{0.004}{10.9091}{mi}{66}\RG{95.806}{0.004}{10.9091}{mi}{118}\RG{104.523}{0.004}{10.9091}{rm}{61}\RG{116.021}{0.004}{10.9091}{rm}{48}\RG{121.493}{0.004}{10.9091}{mi}{58}%
}
\]
又\allowbreak{}\hskip.1em \RB{25.600}{8.311}{2.856}{%
\RG{0.000}{-0.000}{10.9091}{mi}{122}\RG{8.583}{-0.000}{10.9091}{sy}{54}\RG{8.583}{-0.000}{10.9091}{rm}{61}\RG{20.081}{-0.000}{10.9091}{rm}{48}%
}\hskip.1em\allowbreak{} ，所以\allowbreak{}\hskip.1em \RB{16.678}{5.322}{2.606}{%
\RG{0.000}{-0.000}{10.9091}{mi}{117}\RG{6.240}{-0.000}{10.9091}{mi}{59}\RG{11.073}{-0.000}{10.9091}{mi}{118}%
}\hskip.1em\allowbreak{} 中至少一个非零，取它为\allowbreak{}\hskip.1em \RB{7.067}{5.322}{0.620}{%
\RG{0.000}{-0.000}{10.9091}{mi}{11}%
}\hskip.1em\allowbreak{} 即可。\allowbreak{}\hskip.1em \RB{8.366}{7.995}{0.500}{%
\RG{0.000}{-0.000}{10.9091}{am}{3}%
}\hskip.1em\allowbreak{}

方法小结：先对复核向量取Hermite二次型，再比较虚部；关键是半正定性保证\allowbreak{}\hskip.1em \RB{45.660}{8.163}{0.741}{%
\RG{0.000}{-0.001}{10.9091}{mi}{122}\RG{5.554}{3.957}{7.9701}{sy}{3}\RG{10.286}{-0.001}{10.9091}{mi}{66}\RG{19.101}{-0.001}{10.9091}{mi}{122}\RG{28.163}{-0.001}{10.9091}{rm}{61}\RG{40.142}{-0.001}{10.9091}{rm}{48}%
}\hskip.1em\allowbreak{} 推出\allowbreak{}\hskip.1em \RB{34.422}{7.955}{0.736}{%
\RG{0.000}{0.004}{10.9091}{mi}{66}\RG{8.815}{0.004}{10.9091}{mi}{122}\RG{17.405}{0.004}{10.9091}{rm}{61}\RG{28.903}{0.004}{10.9091}{rm}{48}%
}\hskip.1em\allowbreak{} 。最后取实部或虚部得到实核向量。

题目来源：谢启鸿《复旦大学高等代数习题集（2023版）》，第117页。

% Original page 5.
\section*{2022S13}
设实二次型\allowbreak{}\hskip.1em \RB{42.828}{8.685}{3.224}{%
\RG{0.000}{0.003}{10.9091}{mi}{102}\RG{6.515}{0.003}{10.9091}{rm}{40}\RG{10.757}{0.003}{10.9091}{mi}{120}\RG{16.992}{0.003}{10.9091}{rm}{41}\RG{24.650}{0.003}{10.9091}{rm}{61}\RG{36.579}{0.003}{10.9091}{mi}{120}%
}\hskip.1em\allowbreak{} T\allowbreak{}\hskip.1em \RB{14.431}{8.314}{0.617}{%
\RG{0.000}{0.003}{10.9091}{mi}{65}\RG{8.182}{0.003}{10.9091}{mi}{120}%
}\hskip.1em\allowbreak{} ，其中\allowbreak{}\hskip.1em \RB{54.308}{8.685}{3.224}{%
\RG{0.000}{0.003}{10.9091}{mi}{65}\RG{11.612}{0.003}{10.9091}{sy}{50}\RG{22.314}{0.003}{10.9091}{mi}{77}\RG{32.898}{-1.634}{7.9701}{mi}{110}\RG{38.534}{0.003}{10.9091}{rm}{40}\RG{42.776}{0.003}{10.9091}{bb}{82}\RG{50.655}{0.003}{10.9091}{rm}{41}%
}\hskip.1em\allowbreak{} 未必对称，且det\allowbreak{}\hskip.1em \RB{29.044}{8.314}{0.933}{%
\RG{0.000}{0.003}{10.9091}{mi}{65}\RG{11.607}{0.003}{10.9091}{mi}{60}\RG{23.526}{0.003}{10.9091}{rm}{48}%
}\hskip.1em\allowbreak{} 。证明存在实数\allowbreak{}\hskip.1em \RB{45.886}{5.325}{2.602}{%
\RG{0.000}{0.003}{10.9091}{mi}{97}\RG{5.767}{-1.634}{7.9701}{rm}{49}\RG{10.499}{0.003}{10.9091}{mi}{59}\RG{15.332}{0.003}{10.9091}{mi}{58}\RG{20.176}{0.003}{10.9091}{mi}{58}\RG{25.008}{0.003}{10.9091}{mi}{58}\RG{29.852}{0.003}{10.9091}{mi}{59}\RG{34.696}{0.003}{10.9091}{mi}{97}\RG{40.508}{-1.634}{7.9701}{mi}{110}%
}\hskip.1em\allowbreak{} ，使

\[
\RB{83.739}{8.685}{3.224}{%
\RG{0.000}{0.003}{10.9091}{mi}{102}\RG{6.515}{0.003}{10.9091}{rm}{40}\RG{10.757}{0.003}{10.9091}{mi}{97}\RG{21.256}{0.003}{10.9091}{mi}{59}\RG{26.089}{0.003}{10.9091}{mi}{58}\RG{30.933}{0.003}{10.9091}{mi}{58}\RG{35.776}{0.003}{10.9091}{mi}{58}\RG{40.609}{0.003}{10.9091}{mi}{59}\RG{45.453}{0.003}{10.9091}{mi}{97}\RG{56.901}{0.003}{10.9091}{rm}{41}\RG{64.175}{0.003}{10.9091}{mi}{60}\RG{75.690}{0.003}{10.9091}{rm}{48}\RG{81.145}{0.003}{10.9091}{mi}{58}\RG{16.524}{-1.634}{7.9701}{rm}{49}\RG{51.265}{-1.634}{7.9701}{mi}{110}%
}
\]
证明：考虑实系数多项式

\[
\RB{94.577}{9.291}{3.229}{%
\RG{0.000}{-0.002}{10.9091}{mi}{112}\RG{5.488}{-0.002}{10.9091}{rm}{40}\RG{9.731}{-0.002}{10.9091}{mi}{116}\RG{13.670}{-0.002}{10.9091}{rm}{41}\RG{20.925}{-0.002}{10.9091}{rm}{61}\RT{32.458}{-0.002}{10.9091}{d}\RT{38.524}{-0.002}{10.9091}{e}\RT{43.367}{-0.002}{10.9091}{t}\RG{47.611}{-0.002}{10.9091}{rm}{40}\RG{51.853}{-0.002}{10.9091}{mi}{65}\RG{62.460}{-0.002}{10.9091}{rm}{43}\RG{73.369}{-0.002}{10.9091}{mi}{116}\RG{77.307}{-0.002}{10.9091}{mi}{73}\RG{87.740}{-0.002}{10.9091}{rm}{41}\RG{91.983}{-0.002}{10.9091}{mi}{58}\RG{82.104}{-1.638}{7.9701}{mi}{110}%
}
\]
它的最高次项为\allowbreak{}\hskip.1em \RB{9.318}{7.971}{0.623}{%
\RG{0.000}{-0.003}{10.9091}{mi}{116}\RG{3.940}{3.956}{7.9701}{mi}{110}%
}\hskip.1em\allowbreak{} ，且

\[
\RB{188.640}{9.295}{8.047}{%
\RG{0.000}{0.002}{10.9091}{mi}{112}\RG{5.488}{0.002}{10.9091}{rm}{40}\RG{9.721}{0.002}{10.9091}{rm}{48}\RG{15.175}{0.002}{10.9091}{rm}{41}\RG{22.430}{0.002}{10.9091}{rm}{61}\RT{33.973}{0.002}{10.9091}{d}\RT{40.039}{0.002}{10.9091}{e}\RT{44.882}{0.002}{10.9091}{t}\RG{50.944}{0.002}{10.9091}{mi}{65}\RG{62.148}{0.002}{10.9091}{mi}{60}\RG{73.672}{0.002}{10.9091}{rm}{48}\RG{79.127}{0.002}{10.9091}{mi}{59}\RT{111.508}{0.002}{10.9091}{l}\RT{114.541}{0.002}{10.9091}{i}\RT{117.573}{0.002}{10.9091}{m}\RG{134.193}{0.002}{10.9091}{mi}{112}\RG{139.681}{0.002}{10.9091}{rm}{40}\RG{143.924}{0.002}{10.9091}{mi}{116}\RG{147.863}{0.002}{10.9091}{rm}{41}\RG{155.118}{0.002}{10.9091}{rm}{61}\RG{166.627}{0.002}{10.9091}{rm}{43}\RG{175.136}{0.002}{10.9091}{sy}{49}\RG{186.046}{0.002}{10.9091}{mi}{58}\RG{105.793}{-6.718}{7.9701}{mi}{116}\RG{108.851}{-6.718}{7.9701}{sy}{33}\RG{117.320}{-6.718}{7.9701}{rm}{43}\RG{123.906}{-6.718}{7.9701}{sy}{49}%
}
\]
因此可取\allowbreak{}\hskip.1em \RB{27.954}{7.887}{0.935}{%
\RG{0.000}{0.001}{10.9091}{mi}{84}\RG{10.909}{0.001}{10.9091}{mi}{62}\RG{22.436}{0.001}{10.9091}{rm}{48}%
}\hskip.1em\allowbreak{} 使\allowbreak{}\hskip.1em \RB{41.928}{8.683}{3.226}{%
\RG{0.000}{0.001}{10.9091}{mi}{112}\RG{5.488}{0.001}{10.9091}{rm}{40}\RG{9.731}{0.001}{10.9091}{mi}{84}\RG{17.621}{0.001}{10.9091}{rm}{41}\RG{24.894}{0.001}{10.9091}{mi}{62}\RG{36.409}{0.001}{10.9091}{rm}{48}%
}\hskip.1em\allowbreak{} 。由连续函数的介值定理，存在\allowbreak{}\hskip.1em \RB{48.094}{8.683}{3.226}{%
\RG{0.000}{0.001}{10.9091}{mi}{116}\RG{3.940}{-1.635}{7.9701}{rm}{48}\RG{11.702}{0.001}{10.9091}{sy}{50}\RG{22.005}{0.001}{10.9091}{rm}{40}\RG{26.238}{0.001}{10.9091}{rm}{48}\RG{31.702}{0.001}{10.9091}{mi}{59}\RG{36.535}{0.001}{10.9091}{mi}{84}\RG{44.441}{0.001}{10.9091}{rm}{41}%
}\hskip.1em\allowbreak{} ，满足

\[
\RB{86.851}{9.291}{3.229}{%
\RT{0.000}{-0.002}{10.9091}{d}\RT{6.065}{-0.002}{10.9091}{e}\RT{10.909}{-0.002}{10.9091}{t}\RG{15.152}{-0.002}{10.9091}{rm}{40}\RG{19.395}{-0.002}{10.9091}{mi}{65}\RG{30.001}{-0.002}{10.9091}{rm}{43}\RG{40.910}{-0.002}{10.9091}{mi}{116}\RG{49.582}{-0.002}{10.9091}{mi}{73}\RG{60.013}{-0.002}{10.9091}{rm}{41}\RG{67.268}{-0.002}{10.9091}{rm}{61}\RG{78.777}{-0.002}{10.9091}{rm}{48}\RG{84.257}{-0.002}{10.9091}{mi}{58}\RG{44.849}{-1.639}{7.9701}{rm}{48}\RG{54.377}{-1.639}{7.9701}{mi}{110}%
}
\]
矩阵\allowbreak{}\hskip.1em \RB{40.360}{8.315}{2.300}{%
\RG{0.000}{0.004}{10.9091}{mi}{65}\RG{10.606}{0.004}{10.9091}{rm}{43}\RG{21.515}{0.004}{10.9091}{mi}{116}\RG{25.454}{-1.633}{7.9701}{rm}{48}\RG{30.187}{0.004}{10.9091}{mi}{73}\RG{34.982}{-1.633}{7.9701}{mi}{110}%
}\hskip.1em\allowbreak{} 为实奇异矩阵，故齐次实线性方程组

\[
\RB{75.385}{8.682}{3.227}{%
\RG{0.000}{-0.000}{10.9091}{rm}{40}\RG{4.242}{-0.000}{10.9091}{mi}{65}\RG{14.848}{-0.000}{10.9091}{rm}{43}\RG{25.757}{-0.000}{10.9091}{mi}{116}\RG{34.429}{-0.000}{10.9091}{mi}{73}\RG{44.861}{-0.000}{10.9091}{rm}{41}\RG{49.103}{-0.000}{10.9091}{mi}{120}\RG{58.368}{-0.000}{10.9091}{rm}{61}\RG{69.867}{-0.000}{10.9091}{rm}{48}\RG{29.697}{-1.637}{7.9701}{rm}{48}\RG{39.225}{-1.637}{7.9701}{mi}{110}%
}
\]
存在非零实解\allowbreak{}\hskip.1em \RB{74.820}{8.678}{3.231}{%
\RG{0.000}{-0.004}{10.9091}{mi}{120}\RG{9.265}{-0.004}{10.9091}{rm}{61}\RG{20.763}{-0.004}{10.9091}{rm}{40}\RG{25.023}{-0.004}{10.9091}{mi}{97}\RG{30.789}{-1.641}{7.9701}{rm}{49}\RG{35.522}{-0.004}{10.9091}{mi}{59}\RG{40.355}{-0.004}{10.9091}{mi}{58}\RG{45.199}{-0.004}{10.9091}{mi}{58}\RG{50.032}{-0.004}{10.9091}{mi}{58}\RG{54.875}{-0.004}{10.9091}{mi}{59}\RG{59.719}{-0.004}{10.9091}{mi}{97}\RG{65.531}{-1.641}{7.9701}{mi}{110}\RG{71.168}{-0.004}{10.9091}{rm}{41}%
}\hskip.1em\allowbreak{} T。于是\allowbreak{}\hskip.1em \RB{52.368}{8.307}{2.308}{%
\RG{0.000}{-0.004}{10.9091}{mi}{65}\RG{8.182}{-0.004}{10.9091}{mi}{120}\RG{17.446}{-0.004}{10.9091}{rm}{61}\RG{28.962}{-0.004}{10.9091}{sy}{0}\RG{37.446}{-0.004}{10.9091}{mi}{116}\RG{41.386}{-1.641}{7.9701}{rm}{48}\RG{46.118}{-0.004}{10.9091}{mi}{120}%
}\hskip.1em\allowbreak{} ，从而

\[
\RB{237.272}{17.656}{15.209}{%
\RG{191.913}{13.641}{7.9701}{mi}{110}\RG{186.603}{10.368}{10.9092}{ex}{88}\RT{81.925}{4.509}{7.9701}{T}\RT{140.563}{4.509}{7.9701}{T}\RG{209.946}{4.509}{7.9701}{rm}{50}\RG{0.000}{0.004}{10.9091}{mi}{102}\RG{6.515}{0.004}{10.9091}{rm}{40}\RG{10.757}{0.004}{10.9091}{mi}{97}\RG{21.256}{0.004}{10.9091}{mi}{59}\RG{26.089}{0.004}{10.9091}{mi}{58}\RG{30.933}{0.004}{10.9091}{mi}{58}\RG{35.776}{0.004}{10.9091}{mi}{58}\RG{40.609}{0.004}{10.9091}{mi}{59}\RG{45.453}{0.004}{10.9091}{mi}{97}\RG{56.901}{0.004}{10.9091}{rm}{41}\RG{64.156}{0.004}{10.9091}{rm}{61}\RG{75.690}{0.004}{10.9091}{mi}{120}\RG{88.209}{0.004}{10.9091}{mi}{65}\RG{96.391}{0.004}{10.9091}{mi}{120}\RG{105.656}{0.004}{10.9091}{rm}{61}\RG{117.171}{0.004}{10.9091}{sy}{0}\RG{125.656}{0.004}{10.9091}{mi}{116}\RG{134.328}{0.004}{10.9091}{mi}{120}\RG{146.847}{0.004}{10.9091}{mi}{120}\RG{156.112}{0.004}{10.9091}{rm}{61}\RG{167.628}{0.004}{10.9091}{sy}{0}\RG{176.112}{0.004}{10.9091}{mi}{116}\RG{204.179}{0.004}{10.9091}{mi}{97}\RG{217.708}{0.004}{10.9091}{mi}{60}\RG{229.223}{0.004}{10.9091}{rm}{48}\RG{234.678}{0.004}{10.9091}{mi}{58}\RG{16.524}{-1.632}{7.9701}{rm}{49}\RG{51.265}{-1.632}{7.9701}{mi}{110}\RG{129.596}{-1.632}{7.9701}{rm}{48}\RG{180.052}{-1.632}{7.9701}{rm}{48}\RG{209.946}{-2.692}{7.9701}{mi}{106}\RG{187.130}{-13.083}{7.9701}{mi}{106}\RG{191.014}{-13.083}{7.9701}{rm}{61}\RG{197.597}{-13.083}{7.9701}{rm}{49}%
}
\]
这就证明了所求结论。\allowbreak{}\hskip.1em \RB{8.366}{7.990}{0.505}{%
\RG{0.000}{-0.005}{10.9091}{am}{3}%
}\hskip.1em\allowbreak{}

这一论证同时说明：行列式为负的实方阵，必有一个负实特征值及相应的非零实特征向量。全程不需要\allowbreak{}\hskip.1em \RB{8.366}{8.307}{0.504}{%
\RG{0.000}{-0.004}{10.9091}{mi}{65}%
}\hskip.1em\allowbreak{} 对称，也不需要\allowbreak{}\hskip.1em \RB{8.366}{8.307}{0.504}{%
\RG{0.000}{-0.004}{10.9091}{mi}{65}%
}\hskip.1em\allowbreak{} 可对角化。

还需注意，虽然二次型只取决于\allowbreak{}\hskip.1em \RB{8.366}{8.313}{0.500}{%
\RG{0.000}{0.002}{10.9091}{mi}{65}%
}\hskip.1em\allowbreak{} 的对称部分，即

\[
\RB{101.317}{18.260}{7.986}{%
\RT{84.894}{11.336}{7.9701}{T}\RG{55.196}{7.378}{10.9091}{mi}{65}\RG{65.803}{7.378}{10.9091}{rm}{43}\RG{76.712}{7.378}{10.9091}{mi}{65}\RT{6.235}{4.502}{7.9701}{T}\RT{47.716}{4.502}{7.9701}{T}\RG{0.000}{-0.003}{10.9091}{mi}{120}\RG{12.519}{-0.003}{10.9091}{mi}{65}\RG{20.701}{-0.003}{10.9091}{mi}{120}\RG{29.966}{-0.003}{10.9091}{rm}{61}\RG{41.481}{-0.003}{10.9091}{mi}{120}\RG{92.374}{-0.003}{10.9091}{mi}{120}\RG{98.603}{-0.003}{10.9091}{mi}{59}\RG{70.460}{-7.486}{10.9091}{rm}{50}\RR{55.196}{2.725}{35.981}{0.436}%
}
\]
但题设给的是det\allowbreak{}\hskip.1em \RB{28.432}{8.313}{0.934}{%
\RG{0.000}{0.002}{10.9091}{mi}{65}\RG{11.302}{0.002}{10.9091}{mi}{60}\RG{22.914}{0.002}{10.9091}{rm}{48}%
}\hskip.1em\allowbreak{} ，不能未经证明就将它替换成对称部分的行列式条件。上面的证明直接使用\allowbreak{}\hskip.1em \RB{8.366}{8.308}{0.503}{%
\RG{0.000}{-0.003}{10.9091}{mi}{65}%
}\hskip.1em\allowbreak{} ，避免了这一步困难。

方法小结：把行列式的负号转化为多项式在正半轴上的零点，从而找到负实特征值；沿对应实特征向量，二次型立即取负值。

题目来源：谢启鸿《复旦大学高等代数习题集（2023版）》，第117页。

% Original page 6.
\section*{2022S14}
设\allowbreak{}\hskip.1em \RB{21.762}{8.314}{2.602}{%
\RG{0.000}{0.003}{10.9091}{mi}{65}\RG{8.182}{0.003}{10.9091}{mi}{59}\RG{13.015}{0.003}{10.9091}{mi}{66}%
}\hskip.1em\allowbreak{} 均为\allowbreak{}\hskip.1em \RB{6.729}{5.325}{0.617}{%
\RG{0.000}{0.003}{10.9091}{mi}{110}%
}\hskip.1em\allowbreak{} 阶半正定实对称阵。证明

\[
\RB{156.957}{15.142}{8.107}{%
\RG{0.547}{7.376}{10.9091}{rm}{49}\RG{91.755}{4.500}{7.9701}{rm}{49}\RG{95.987}{4.500}{7.9701}{rm}{47}\RG{100.223}{4.500}{7.9701}{mi}{110}\RG{140.137}{4.500}{7.9701}{rm}{49}\RG{144.369}{4.500}{7.9701}{rm}{47}\RG{148.605}{4.500}{7.9701}{mi}{110}\RT{9.562}{-0.004}{10.9091}{t}\RT{13.806}{-0.004}{10.9091}{r}\RG{18.082}{-0.004}{10.9091}{rm}{40}\RG{22.324}{-0.004}{10.9091}{mi}{65}\RG{30.506}{-0.004}{10.9091}{mi}{66}\RG{39.328}{-0.004}{10.9091}{rm}{41}\RG{46.601}{-0.004}{10.9091}{sy}{21}\RG{58.116}{-0.004}{10.9091}{rm}{40}\RT{62.359}{-0.004}{10.9091}{d}\RT{68.424}{-0.004}{10.9091}{e}\RT{73.268}{-0.004}{10.9091}{t}\RG{79.330}{-0.004}{10.9091}{mi}{65}\RG{87.512}{-0.004}{10.9091}{rm}{41}\RG{105.859}{-0.004}{10.9091}{rm}{40}\RT{110.102}{-0.004}{10.9091}{d}\RT{116.167}{-0.004}{10.9091}{e}\RT{121.011}{-0.004}{10.9091}{t}\RG{127.073}{-0.004}{10.9091}{mi}{66}\RG{135.895}{-0.004}{10.9091}{rm}{41}\RG{154.242}{-0.004}{10.9091}{mi}{59}\RG{0.000}{-7.487}{10.9091}{mi}{110}\RR{0.000}{2.723}{6.548}{0.436}%
}
\]
并求等号成立的充要条件。证明：由正交对角化，\allowbreak{}\hskip.1em \RB{8.366}{8.309}{0.502}{%
\RG{0.000}{-0.002}{10.9091}{mi}{65}%
}\hskip.1em\allowbreak{} 存在半正定平方根\allowbreak{}\hskip.1em \RB{20.948}{10.434}{0.502}{%
\RG{0.000}{-0.002}{10.9091}{mi}{65}\RG{8.182}{3.956}{7.9701}{rm}{49}\RG{12.414}{3.956}{7.9701}{rm}{47}\RG{16.646}{3.956}{7.9701}{rm}{50}%
}\hskip.1em\allowbreak{} 。令

\[
\RB{77.305}{6.476}{5.246}{%
\RG{31.304}{-0.002}{7.9701}{rm}{49}\RG{35.536}{-0.002}{7.9701}{rm}{47}\RG{39.769}{-0.002}{7.9701}{rm}{50}\RG{61.510}{-0.002}{7.9701}{rm}{49}\RG{65.742}{-0.002}{7.9701}{rm}{47}\RG{69.974}{-0.002}{7.9701}{rm}{50}\RG{0.000}{-4.506}{10.9091}{mi}{67}\RG{11.607}{-4.506}{10.9091}{rm}{61}\RG{23.123}{-4.506}{10.9091}{mi}{65}\RG{44.505}{-4.506}{10.9091}{mi}{66}\RG{53.320}{-4.506}{10.9091}{mi}{65}\RG{74.711}{-4.506}{10.9091}{mi}{58}%
}
\]
对任意\allowbreak{}\hskip.1em \RB{35.390}{7.974}{0.936}{%
\RG{0.000}{-0.000}{10.9091}{mi}{120}\RG{10.548}{-0.000}{10.9091}{sy}{50}\RG{22.133}{-0.000}{10.9091}{bb}{82}\RG{30.012}{3.959}{7.9701}{mi}{110}%
}\hskip.1em\allowbreak{} ，有\allowbreak{}\hskip.1em \RB{6.249}{5.322}{0.620}{%
\RG{0.000}{-0.000}{10.9091}{mi}{120}%
}\hskip.1em\allowbreak{} T\allowbreak{}\hskip.1em \RB{67.435}{10.437}{3.227}{%
\RG{0.000}{-0.000}{10.9091}{mi}{67}\RG{8.564}{-0.000}{10.9091}{mi}{120}\RG{19.124}{-0.000}{10.9091}{rm}{61}\RG{31.910}{-0.000}{10.9091}{rm}{40}\RG{36.164}{-0.000}{10.9091}{mi}{65}\RG{44.346}{3.959}{7.9701}{rm}{49}\RG{48.579}{3.959}{7.9701}{rm}{47}\RG{52.811}{3.959}{7.9701}{rm}{50}\RG{57.547}{-0.000}{10.9091}{mi}{120}\RG{63.783}{-0.000}{10.9091}{rm}{41}%
}\hskip.1em\allowbreak{} T\allowbreak{}\hskip.1em \RB{67.553}{10.437}{3.227}{%
\RG{0.000}{-0.000}{10.9091}{mi}{66}\RG{8.822}{-0.000}{10.9091}{rm}{40}\RG{13.064}{-0.000}{10.9091}{mi}{65}\RG{21.246}{3.959}{7.9701}{rm}{49}\RG{25.478}{3.959}{7.9701}{rm}{47}\RG{29.710}{3.959}{7.9701}{rm}{50}\RG{34.447}{-0.000}{10.9091}{mi}{120}\RG{40.681}{-0.000}{10.9091}{rm}{41}\RG{49.237}{-0.000}{10.9091}{sy}{21}\RG{62.035}{-0.000}{10.9091}{rm}{48}%
}\hskip.1em\allowbreak{} ，故\allowbreak{}\hskip.1em \RB{8.791}{8.191}{0.740}{%
\RG{0.000}{-0.000}{10.9091}{mi}{67}%
}\hskip.1em\allowbreak{} 为半正定实对称阵。记其特征值为

\[
\RB{78.748}{9.232}{2.677}{%
\RG{0.000}{0.005}{10.9091}{mi}{21}\RG{11.096}{0.005}{10.9091}{mi}{59}\RG{15.929}{0.005}{10.9091}{mi}{58}\RG{20.773}{0.005}{10.9091}{mi}{58}\RG{25.605}{0.005}{10.9091}{mi}{58}\RG{30.449}{0.005}{10.9091}{mi}{59}\RG{35.293}{0.005}{10.9091}{mi}{21}\RG{50.368}{0.005}{10.9091}{sy}{21}\RG{61.885}{0.005}{10.9091}{rm}{48}\RT{67.339}{0.005}{10.9091}{。}\RG{6.364}{-1.632}{7.9701}{rm}{49}\RG{41.702}{-1.632}{7.9701}{mi}{110}%
}
\]
利用迹的循环性与行列式的乘法性，得到

\[
\RB{145.728}{55.523}{15.214}{%
\RG{93.992}{51.508}{7.9701}{mi}{110}\RG{88.682}{48.235}{10.9092}{ex}{88}\RT{0.000}{37.871}{10.9091}{t}\RT{4.244}{37.871}{10.9091}{r}\RG{8.520}{37.871}{10.9091}{rm}{40}\RG{12.762}{37.871}{10.9091}{mi}{65}\RG{20.944}{37.871}{10.9091}{mi}{66}\RG{29.766}{37.871}{10.9091}{rm}{41}\RG{37.021}{37.871}{10.9091}{rm}{61}\RT{48.554}{37.871}{10.9091}{t}\RT{52.798}{37.871}{10.9091}{r}\RG{57.074}{37.871}{10.9091}{rm}{40}\RG{61.317}{37.871}{10.9091}{mi}{67}\RG{69.895}{37.871}{10.9091}{rm}{41}\RG{77.149}{37.871}{10.9091}{rm}{61}\RG{106.258}{37.871}{10.9091}{mi}{21}\RG{117.004}{37.871}{10.9091}{mi}{59}\RG{112.622}{36.235}{7.9701}{mi}{106}\RG{89.209}{24.784}{7.9701}{mi}{106}\RG{93.093}{24.784}{7.9701}{rm}{61}\RG{99.676}{24.784}{7.9701}{rm}{49}\RG{120.647}{13.636}{7.9701}{mi}{110}\RG{116.247}{10.363}{10.9092}{ex}{89}\RT{8.460}{0.000}{10.9091}{d}\RT{14.526}{0.000}{10.9091}{e}\RT{19.369}{0.000}{10.9091}{t}\RG{25.431}{0.000}{10.9091}{mi}{67}\RG{37.039}{0.000}{10.9091}{rm}{61}\RT{48.554}{0.000}{10.9091}{d}\RT{54.620}{0.000}{10.9091}{e}\RT{59.464}{0.000}{10.9091}{t}\RG{65.526}{0.000}{10.9091}{mi}{65}\RT{75.526}{0.000}{10.9091}{d}\RT{81.591}{0.000}{10.9091}{e}\RT{86.435}{0.000}{10.9091}{t}\RG{92.497}{0.000}{10.9091}{mi}{66}\RG{104.349}{0.000}{10.9091}{rm}{61}\RG{132.387}{0.000}{10.9091}{mi}{21}\RG{143.133}{0.000}{10.9091}{mi}{58}\RG{138.750}{-1.637}{7.9701}{mi}{106}\RG{115.864}{-13.088}{7.9701}{mi}{106}\RG{119.748}{-13.088}{7.9701}{rm}{61}\RG{126.331}{-13.088}{7.9701}{rm}{49}%
}
\]
对这些非负数应用算术—几何平均不等式，便有

\[
\RB{312.858}{15.454}{8.100}{%
\RG{0.547}{7.383}{10.9091}{rm}{49}\RG{59.311}{7.383}{10.9091}{mi}{21}\RG{72.832}{7.383}{10.9091}{rm}{43}\RG{83.741}{7.383}{10.9091}{sy}{1}\RG{88.573}{7.383}{10.9091}{sy}{1}\RG{93.417}{7.383}{10.9091}{sy}{1}\RG{98.893}{7.383}{10.9091}{rm}{43}\RG{109.802}{7.383}{10.9091}{mi}{21}\RG{65.676}{5.747}{7.9701}{rm}{49}\RG{116.166}{5.747}{7.9701}{mi}{110}\RG{185.488}{4.507}{7.9701}{rm}{49}\RG{189.720}{4.507}{7.9701}{rm}{47}\RG{193.957}{4.507}{7.9701}{mi}{110}\RG{247.776}{4.507}{7.9701}{rm}{49}\RG{252.008}{4.507}{7.9701}{rm}{47}\RG{256.244}{4.507}{7.9701}{mi}{110}\RG{296.158}{4.507}{7.9701}{rm}{49}\RG{300.391}{4.507}{7.9701}{rm}{47}\RG{304.628}{4.507}{7.9701}{mi}{110}\RT{9.562}{0.003}{10.9091}{t}\RT{13.806}{0.003}{10.9091}{r}\RG{18.082}{0.003}{10.9091}{rm}{40}\RG{22.324}{0.003}{10.9091}{mi}{65}\RG{30.506}{0.003}{10.9091}{mi}{66}\RG{39.328}{0.003}{10.9091}{rm}{41}\RG{46.583}{0.003}{10.9091}{rm}{61}\RG{126.028}{0.003}{10.9091}{sy}{21}\RG{137.543}{0.003}{10.9091}{rm}{40}\RG{141.785}{0.003}{10.9091}{mi}{21}\RG{154.699}{0.003}{10.9091}{sy}{1}\RG{159.532}{0.003}{10.9091}{sy}{1}\RG{164.376}{0.003}{10.9091}{sy}{1}\RG{169.245}{0.003}{10.9091}{mi}{21}\RG{181.245}{0.003}{10.9091}{rm}{41}\RG{202.623}{0.003}{10.9091}{rm}{61}\RG{214.121}{0.003}{10.9091}{rm}{40}\RT{218.381}{0.003}{10.9091}{d}\RT{224.446}{0.003}{10.9091}{e}\RT{229.290}{0.003}{10.9091}{t}\RG{235.352}{0.003}{10.9091}{mi}{65}\RG{243.533}{0.003}{10.9091}{rm}{41}\RG{261.881}{0.003}{10.9091}{rm}{40}\RT{266.123}{0.003}{10.9091}{d}\RT{272.189}{0.003}{10.9091}{e}\RT{277.032}{0.003}{10.9091}{t}\RG{283.094}{0.003}{10.9091}{mi}{66}\RG{291.916}{0.003}{10.9091}{rm}{41}\RG{310.264}{0.003}{10.9091}{mi}{58}\RG{148.149}{-1.633}{7.9701}{rm}{49}\RG{175.608}{-1.633}{7.9701}{mi}{110}\RG{0.000}{-7.480}{10.9091}{mi}{110}\RG{87.283}{-7.480}{10.9091}{mi}{110}\RR{0.000}{2.730}{6.548}{0.436}\RR{59.311}{2.730}{62.491}{0.436}%
}
\]
等号成立当且仅当\allowbreak{}\hskip.1em \RB{106.116}{8.076}{2.211}{%
\RG{0.000}{0.005}{10.9091}{mi}{21}\RG{6.364}{-1.631}{7.9701}{rm}{49}\RG{14.361}{0.005}{10.9091}{rm}{61}\RG{26.109}{0.005}{10.9091}{sy}{1}\RG{30.942}{0.005}{10.9091}{sy}{1}\RG{35.786}{0.005}{10.9091}{sy}{1}\RG{42.101}{0.005}{10.9091}{rm}{61}\RG{53.850}{0.005}{10.9091}{mi}{21}\RG{60.214}{-1.631}{7.9701}{mi}{110}\RG{69.115}{0.005}{10.9091}{rm}{61}\RG{80.864}{0.005}{10.9091}{mi}{99}\RG{88.849}{0.005}{10.9091}{sy}{21}\RG{100.598}{0.005}{10.9091}{rm}{48}%
}\hskip.1em\allowbreak{} 。由于\allowbreak{}\hskip.1em \RB{8.791}{8.196}{0.735}{%
\RG{0.000}{0.005}{10.9091}{mi}{67}%
}\hskip.1em\allowbreak{} 为实对称阵，这又等价于\allowbreak{}\hskip.1em \RB{38.485}{8.196}{2.211}{%
\RG{0.000}{0.005}{10.9091}{mi}{67}\RG{11.841}{0.005}{10.9091}{rm}{61}\RG{23.590}{0.005}{10.9091}{mi}{99}\RG{28.303}{0.005}{10.9091}{mi}{73}\RG{33.107}{-1.631}{7.9701}{mi}{110}%
}\hskip.1em\allowbreak{} 。分两种情形讨论。

若\allowbreak{}\hskip.1em \RB{24.784}{7.762}{0.940}{%
\RG{0.000}{-0.004}{10.9091}{mi}{99}\RG{7.735}{-0.004}{10.9091}{mi}{62}\RG{19.266}{-0.004}{10.9091}{rm}{48}%
}\hskip.1em\allowbreak{} ，则\allowbreak{}\hskip.1em \RB{8.791}{8.187}{0.744}{%
\RG{0.000}{-0.004}{10.9091}{mi}{67}%
}\hskip.1em\allowbreak{} 可逆，因而\allowbreak{}\hskip.1em \RB{8.366}{8.307}{0.504}{%
\RG{0.000}{-0.004}{10.9091}{mi}{65}%
}\hskip.1em\allowbreak{} 可逆，并且

\[
\RB{202.614}{10.984}{2.603}{%
\RG{31.549}{4.506}{7.9701}{sy}{0}\RG{38.135}{4.506}{7.9701}{rm}{49}\RG{42.367}{4.506}{7.9701}{rm}{47}\RG{46.600}{4.506}{7.9701}{rm}{50}\RG{68.096}{4.506}{7.9701}{sy}{0}\RG{74.683}{4.506}{7.9701}{rm}{49}\RG{78.915}{4.506}{7.9701}{rm}{47}\RG{83.147}{4.506}{7.9701}{rm}{50}\RG{115.331}{4.506}{7.9701}{sy}{0}\RG{121.918}{4.506}{7.9701}{rm}{49}\RG{0.000}{0.002}{10.9091}{mi}{66}\RG{11.852}{0.002}{10.9091}{rm}{61}\RG{23.367}{0.002}{10.9091}{mi}{65}\RG{51.336}{0.002}{10.9091}{mi}{67}\RG{59.900}{0.002}{10.9091}{mi}{65}\RG{90.914}{0.002}{10.9091}{rm}{61}\RG{102.429}{0.002}{10.9091}{mi}{99}\RG{107.142}{0.002}{10.9091}{mi}{65}\RG{126.650}{0.002}{10.9091}{mi}{59}\RG{153.301}{0.002}{10.9091}{mi}{65}\RG{161.483}{0.002}{10.9091}{mi}{66}\RG{173.351}{0.002}{10.9091}{rm}{61}\RG{184.867}{0.002}{10.9091}{mi}{99}\RG{189.579}{0.002}{10.9091}{mi}{73}\RG{200.020}{0.002}{10.9091}{mi}{58}\RG{194.384}{-1.634}{7.9701}{mi}{110}%
}
\]
若\allowbreak{}\hskip.1em \RB{24.768}{7.765}{0.741}{%
\RG{0.000}{-0.001}{10.9091}{mi}{99}\RG{7.751}{-0.001}{10.9091}{rm}{61}\RG{19.249}{-0.001}{10.9091}{rm}{48}%
}\hskip.1em\allowbreak{} ，取\allowbreak{}\hskip.1em \RB{8.747}{7.950}{0.501}{%
\RG{0.000}{-0.001}{10.9091}{mi}{66}%
}\hskip.1em\allowbreak{} 的半正定平方根，则

\[
\RB{155.783}{6.924}{7.732}{%
\RG{56.187}{-0.000}{7.9701}{rm}{49}\RG{60.420}{-0.000}{7.9701}{rm}{47}\RG{64.652}{-0.000}{7.9701}{rm}{50}\RG{77.571}{-0.000}{7.9701}{rm}{49}\RG{81.803}{-0.000}{7.9701}{rm}{47}\RG{86.035}{-0.000}{7.9701}{rm}{50}\RT{95.014}{-0.000}{7.9701}{T}\RG{114.363}{-0.000}{7.9701}{rm}{49}\RG{118.595}{-0.000}{7.9701}{rm}{47}\RG{122.827}{-0.000}{7.9701}{rm}{50}\RG{135.745}{-0.000}{7.9701}{rm}{49}\RG{139.977}{-0.000}{7.9701}{rm}{47}\RG{144.209}{-0.000}{7.9701}{rm}{50}\RG{0.000}{-4.505}{10.9091}{rm}{48}\RG{8.476}{-4.505}{10.9091}{rm}{61}\RG{20.000}{-4.505}{10.9091}{mi}{67}\RG{31.608}{-4.505}{10.9091}{rm}{61}\RG{43.107}{-4.505}{10.9091}{rm}{40}\RG{47.365}{-4.505}{10.9091}{mi}{66}\RG{69.389}{-4.505}{10.9091}{mi}{65}\RG{90.771}{-4.505}{10.9091}{rm}{41}\RG{101.298}{-4.505}{10.9091}{rm}{40}\RG{105.541}{-4.505}{10.9091}{mi}{66}\RG{127.563}{-4.505}{10.9091}{mi}{65}\RG{148.946}{-4.505}{10.9091}{rm}{41}\RG{153.188}{-4.505}{10.9091}{mi}{58}%
}
\]
矩阵\allowbreak{}\hskip.1em \RB{11.889}{7.953}{0.500}{%
\RG{0.000}{0.002}{10.9091}{mi}{77}%
}\hskip.1em\allowbreak{} T\allowbreak{}\hskip.1em \RB{31.820}{7.953}{0.738}{%
\RG{0.000}{0.002}{10.9091}{mi}{77}\RG{14.803}{0.002}{10.9091}{rm}{61}\RG{26.301}{0.002}{10.9091}{rm}{48}%
}\hskip.1em\allowbreak{} 蕴含\allowbreak{}\hskip.1em \RB{31.820}{7.953}{0.738}{%
\RG{0.000}{0.002}{10.9091}{mi}{77}\RG{14.803}{0.002}{10.9091}{rm}{61}\RG{26.301}{0.002}{10.9091}{rm}{48}%
}\hskip.1em\allowbreak{} ，因为其每个对角元都是\allowbreak{}\hskip.1em \RB{11.889}{7.953}{0.500}{%
\RG{0.000}{0.002}{10.9091}{mi}{77}%
}\hskip.1em\allowbreak{} 相应列向量的模平方。因此\allowbreak{}\hskip.1em \RB{63.452}{10.438}{0.738}{%
\RG{0.000}{0.002}{10.9091}{mi}{66}\RG{8.822}{3.960}{7.9701}{rm}{49}\RG{13.054}{3.960}{7.9701}{rm}{47}\RG{17.286}{3.960}{7.9701}{rm}{50}\RG{22.022}{0.002}{10.9091}{mi}{65}\RG{30.204}{3.960}{7.9701}{rm}{49}\RG{34.436}{3.960}{7.9701}{rm}{47}\RG{38.669}{3.960}{7.9701}{rm}{50}\RG{46.435}{0.002}{10.9091}{rm}{61}\RG{57.934}{0.002}{10.9091}{rm}{48}%
}\hskip.1em\allowbreak{} ，继而

\[
\RB{259.430}{11.430}{3.226}{%
\RG{40.371}{4.506}{7.9701}{rm}{49}\RG{44.603}{4.506}{7.9701}{rm}{47}\RG{48.836}{4.506}{7.9701}{rm}{50}\RG{66.637}{4.506}{7.9701}{rm}{49}\RG{70.869}{4.506}{7.9701}{rm}{47}\RG{75.101}{4.506}{7.9701}{rm}{50}\RG{88.020}{4.506}{7.9701}{rm}{49}\RG{92.252}{4.506}{7.9701}{rm}{47}\RG{96.484}{4.506}{7.9701}{rm}{50}\RG{113.645}{4.506}{7.9701}{rm}{49}\RG{117.877}{4.506}{7.9701}{rm}{47}\RG{122.109}{4.506}{7.9701}{rm}{50}\RT{230.551}{4.506}{7.9701}{T}\RG{0.000}{0.002}{10.9091}{mi}{66}\RG{8.815}{0.002}{10.9091}{mi}{65}\RG{20.034}{0.001}{10.9091}{rm}{61}\RG{31.549}{0.002}{10.9091}{mi}{66}\RG{53.572}{0.001}{10.9091}{rm}{40}\RG{57.814}{0.002}{10.9091}{mi}{66}\RG{79.838}{0.001}{10.9091}{mi}{65}\RG{101.221}{0.001}{10.9091}{rm}{41}\RG{105.463}{0.002}{10.9091}{mi}{65}\RG{129.876}{0.001}{10.9091}{rm}{61}\RG{141.374}{0.001}{10.9091}{rm}{48}\RG{146.846}{0.002}{10.9091}{mi}{59}\RG{173.497}{0.002}{10.9091}{mi}{65}\RG{181.679}{0.002}{10.9091}{mi}{66}\RG{193.547}{0.001}{10.9091}{rm}{61}\RG{205.045}{0.001}{10.9091}{rm}{40}\RG{209.305}{0.002}{10.9091}{mi}{66}\RG{218.119}{0.002}{10.9091}{mi}{65}\RG{226.309}{0.001}{10.9091}{rm}{41}\RG{239.866}{0.001}{10.9091}{rm}{61}\RG{251.364}{0.001}{10.9091}{rm}{48}\RG{256.835}{0.002}{10.9091}{mi}{58}%
}
\]
反之，若\allowbreak{}\hskip.1em \RB{47.102}{8.313}{2.214}{%
\RG{0.000}{0.002}{10.9091}{mi}{65}\RG{8.182}{0.002}{10.9091}{mi}{66}\RG{20.363}{0.002}{10.9091}{rm}{61}\RG{32.208}{0.002}{10.9091}{mi}{99}\RG{36.921}{0.002}{10.9091}{mi}{73}\RG{41.724}{-1.634}{7.9701}{mi}{110}%
}\hskip.1em\allowbreak{} ，其中\allowbreak{}\hskip.1em \RB{25.442}{7.768}{1.993}{%
\RG{0.000}{0.002}{10.9091}{mi}{99}\RG{8.080}{0.002}{10.9091}{sy}{21}\RG{19.924}{0.002}{10.9091}{rm}{48}%
}\hskip.1em\allowbreak{} ，则tr\allowbreak{}\hskip.1em \RB{52.429}{8.684}{3.225}{%
\RG{0.000}{0.002}{10.9091}{rm}{40}\RG{4.242}{0.002}{10.9091}{mi}{65}\RG{12.424}{0.002}{10.9091}{mi}{66}\RG{21.246}{0.002}{10.9091}{rm}{41}\RG{28.828}{0.002}{10.9091}{rm}{61}\RG{40.692}{0.002}{10.9091}{mi}{110}\RG{47.238}{0.002}{10.9091}{mi}{99}%
}\hskip.1em\allowbreak{} ，det\allowbreak{}\hskip.1em \RB{8.366}{8.313}{0.500}{%
\RG{0.000}{0.002}{10.9091}{mi}{65}%
}\hskip.1em\allowbreak{} det\allowbreak{}\hskip.1em \RB{34.125}{7.976}{0.618}{%
\RG{0.000}{0.002}{10.9091}{mi}{66}\RG{12.182}{0.002}{10.9091}{rm}{61}\RG{24.026}{0.002}{10.9091}{mi}{99}\RG{28.747}{3.961}{7.9701}{mi}{110}%
}\hskip.1em\allowbreak{} ，原不等式两端都等于\allowbreak{}\hskip.1em \RB{5.191}{5.319}{0.623}{%
\RG{0.000}{-0.003}{10.9091}{mi}{99}%
}\hskip.1em\allowbreak{} ，故等号成立。

综上，等号成立的充要条件为

\allowbreak{}\hskip.1em \RB{46.443}{8.310}{2.217}{%
\RG{0.000}{-0.001}{10.9091}{mi}{65}\RG{8.182}{-0.001}{10.9091}{mi}{66}\RG{20.034}{-0.001}{10.9091}{rm}{61}\RG{31.549}{-0.001}{10.9091}{mi}{99}\RG{36.262}{-0.001}{10.9091}{mi}{73}\RG{41.065}{-1.637}{7.9701}{mi}{110}%
}\hskip.1em\allowbreak{} 对某个\allowbreak{}\hskip.1em \RB{24.784}{7.765}{1.996}{%
\RG{0.000}{-0.001}{10.9091}{mi}{99}\RG{7.751}{-0.001}{10.9091}{sy}{21}\RG{19.266}{-0.001}{10.9091}{rm}{48}%
}\hskip.1em\allowbreak{} 成立\allowbreak{}\hskip.1em \RB{2.595}{1.655}{0.501}{%
\RG{0.000}{-0.001}{10.9091}{mi}{58}%
}\hskip.1em\allowbreak{}

等价地：或者\allowbreak{}\hskip.1em \RB{21.762}{8.306}{2.610}{%
\RG{0.000}{-0.005}{10.9091}{mi}{65}\RG{8.182}{-0.005}{10.9091}{mi}{59}\RG{13.015}{-0.005}{10.9091}{mi}{66}%
}\hskip.1em\allowbreak{} 均正定且\allowbreak{}\hskip.1em \RB{46.920}{9.754}{0.625}{%
\RG{0.000}{-0.005}{10.9091}{mi}{66}\RG{11.852}{-0.005}{10.9091}{rm}{61}\RG{23.368}{-0.005}{10.9091}{mi}{99}\RG{28.081}{-0.005}{10.9091}{mi}{65}\RG{36.270}{3.954}{7.9701}{sy}{0}\RG{42.857}{3.954}{7.9701}{rm}{49}%
}\hskip.1em\allowbreak{} （\allowbreak{}\hskip.1em \RB{24.785}{7.761}{0.941}{%
\RG{0.000}{-0.005}{10.9091}{mi}{99}\RG{7.735}{-0.005}{10.9091}{mi}{62}\RG{19.267}{-0.005}{10.9091}{rm}{48}%
}\hskip.1em\allowbreak{} ），或者\allowbreak{}\hskip.1em \RB{37.051}{8.306}{0.745}{%
\RG{0.000}{-0.005}{10.9091}{mi}{65}\RG{8.182}{-0.005}{10.9091}{mi}{66}\RG{20.034}{-0.005}{10.9091}{rm}{61}\RG{31.532}{-0.005}{10.9091}{rm}{48}%
}\hskip.1em\allowbreak{} 。\allowbreak{}\hskip.1em \RB{8.366}{7.990}{0.505}{%
\RG{0.000}{-0.005}{10.9091}{am}{3}%
}\hskip.1em\allowbreak{}

方法小结：用\allowbreak{}\hskip.1em \RB{51.152}{10.441}{0.500}{%
\RG{0.000}{0.004}{10.9091}{mi}{65}\RG{8.182}{3.963}{7.9701}{rm}{49}\RG{12.414}{3.963}{7.9701}{rm}{47}\RG{16.646}{3.963}{7.9701}{rm}{50}\RG{21.383}{0.004}{10.9091}{mi}{66}\RG{30.198}{0.004}{10.9091}{mi}{65}\RG{38.386}{3.963}{7.9701}{rm}{49}\RG{42.618}{3.963}{7.9701}{rm}{47}\RG{46.851}{3.963}{7.9701}{rm}{50}%
}\hskip.1em\allowbreak{} 将乘积转化为半正定对称阵，再对特征值用均值不等式；等号条件须同时处理正定与奇异两种情形。

题目来源：谢启鸿《复旦大学高等代数习题集（2023版）》，第117页。

% Original page 7.
\section*{2022S15}
证明：\allowbreak{}\hskip.1em \RB{6.729}{5.325}{0.617}{%
\RG{0.000}{0.003}{10.9091}{mi}{110}%
}\hskip.1em\allowbreak{} 维欧氏空间\allowbreak{}\hskip.1em \RB{8.889}{7.954}{0.737}{%
\RG{0.000}{0.003}{10.9091}{mi}{86}%
}\hskip.1em\allowbreak{} 中，两两夹角大于直角的向量至多有\allowbreak{}\hskip.1em \RB{24.248}{7.769}{1.402}{%
\RG{0.000}{0.003}{10.9091}{mi}{110}\RG{8.628}{0.003}{10.9091}{rm}{43}\RG{19.177}{0.003}{10.9091}{rm}{49}%
}\hskip.1em\allowbreak{} 个，并举例说明可以取到\allowbreak{}\hskip.1em \RB{24.247}{7.769}{1.402}{%
\RG{0.000}{0.003}{10.9091}{mi}{110}\RG{8.627}{0.003}{10.9091}{rm}{43}\RG{19.176}{0.003}{10.9091}{rm}{49}%
}\hskip.1em\allowbreak{}个这样的向量。

证明：设\allowbreak{}\hskip.1em \RB{47.279}{5.326}{2.601}{%
\RG{0.000}{0.004}{10.9091}{mi}{118}\RG{5.288}{-1.632}{7.9701}{rm}{49}\RG{10.020}{0.004}{10.9091}{mi}{59}\RG{14.853}{0.004}{10.9091}{mi}{58}\RG{19.697}{0.004}{10.9091}{mi}{58}\RG{24.540}{0.004}{10.9091}{mi}{58}\RG{29.373}{0.004}{10.9091}{mi}{59}\RG{34.217}{0.004}{10.9091}{mi}{118}\RG{39.550}{-1.632}{7.9701}{mi}{109}%
}\hskip.1em\allowbreak{} 是满足条件的非零向量，则

\[
\RB{107.991}{8.682}{3.762}{%
\RG{0.000}{0.000}{10.9091}{sy}{104}\RG{4.242}{0.000}{10.9091}{mi}{118}\RG{12.911}{0.000}{10.9091}{mi}{59}\RG{17.744}{0.000}{10.9091}{mi}{118}\RG{27.430}{0.000}{10.9091}{sy}{105}\RG{34.702}{0.000}{10.9091}{mi}{60}\RG{46.218}{0.000}{10.9091}{rm}{48}\RG{73.480}{0.000}{10.9091}{rm}{40}\RG{77.734}{0.000}{10.9091}{mi}{105}\RG{84.522}{0.000}{10.9091}{sy}{54}\RG{84.522}{0.000}{10.9091}{rm}{61}\RG{96.037}{0.000}{10.9091}{mi}{106}\RG{101.154}{0.000}{10.9091}{rm}{41}\RG{105.397}{0.000}{10.9091}{mi}{58}\RG{9.530}{-1.636}{7.9701}{mi}{105}\RG{23.048}{-1.636}{7.9701}{mi}{106}%
}
\]
先证\allowbreak{}\hskip.1em \RB{49.065}{7.762}{1.999}{%
\RG{0.000}{-0.004}{10.9091}{mi}{109}\RG{12.608}{-0.004}{10.9091}{sy}{20}\RG{24.123}{-0.004}{10.9091}{mi}{110}\RG{33.096}{-0.004}{10.9091}{rm}{43}\RG{43.995}{-0.004}{10.9091}{rm}{49}%
}\hskip.1em\allowbreak{} 。若\allowbreak{}\hskip.1em \RB{49.393}{7.762}{1.999}{%
\RG{0.000}{-0.004}{10.9091}{mi}{109}\RG{12.609}{-0.004}{10.9091}{sy}{21}\RG{24.124}{-0.004}{10.9091}{mi}{110}\RG{33.096}{-0.004}{10.9091}{rm}{43}\RG{43.995}{-0.004}{10.9091}{rm}{50}%
}\hskip.1em\allowbreak{} ，考虑线性映射

\[
\RB{280.740}{19.596}{14.121}{%
\RG{190.548}{18.659}{10.9092}{ex}{32}\RG{267.690}{18.659}{10.9092}{ex}{33}\RG{203.318}{13.641}{7.9701}{mi}{109}\RG{246.146}{13.641}{7.9701}{mi}{109}\RG{199.184}{10.368}{10.9092}{ex}{88}\RG{242.012}{10.368}{10.9092}{ex}{88}\RG{24.288}{4.509}{7.9701}{mi}{109}\RG{0.000}{0.004}{10.9091}{mi}{70}\RG{9.743}{0.004}{10.9091}{rm}{58}\RG{16.409}{0.004}{10.9091}{bb}{82}\RG{35.307}{0.004}{10.9091}{sy}{0}\RG{41.973}{0.004}{10.9091}{sy}{33}\RG{55.913}{0.004}{10.9091}{mi}{86}\RG{67.126}{0.004}{10.9091}{sy}{8}\RG{78.035}{0.004}{10.9091}{bb}{82}\RG{85.914}{0.004}{10.9091}{mi}{59}\RG{112.565}{0.004}{10.9091}{mi}{70}\RG{121.111}{0.004}{10.9091}{rm}{40}\RG{125.354}{0.004}{10.9091}{mi}{99}\RG{134.807}{0.004}{10.9091}{mi}{59}\RG{139.640}{0.004}{10.9091}{mi}{58}\RG{144.484}{0.004}{10.9091}{mi}{58}\RG{149.316}{0.004}{10.9091}{mi}{58}\RG{154.160}{0.004}{10.9091}{mi}{59}\RG{159.004}{0.004}{10.9091}{mi}{99}\RG{171.759}{0.004}{10.9091}{rm}{41}\RG{179.014}{0.004}{10.9091}{rm}{61}\RG{216.760}{0.004}{10.9091}{mi}{99}\RG{224.862}{0.004}{10.9091}{mi}{118}\RG{233.531}{0.004}{10.9091}{mi}{59}\RG{259.588}{0.004}{10.9091}{mi}{99}\RG{278.145}{0.004}{10.9091}{mi}{58}\RG{130.075}{-1.632}{7.9701}{rm}{49}\RG{163.770}{-1.632}{7.9701}{mi}{109}\RG{221.481}{-1.632}{7.9701}{mi}{105}\RG{230.150}{-1.632}{7.9701}{mi}{105}\RG{264.309}{-1.632}{7.9701}{mi}{105}\RG{200.211}{-13.083}{7.9701}{mi}{105}\RG{203.094}{-13.083}{7.9701}{rm}{61}\RG{209.677}{-13.083}{7.9701}{rm}{49}\RG{243.039}{-13.083}{7.9701}{mi}{105}\RG{245.922}{-13.083}{7.9701}{rm}{61}\RG{252.505}{-13.083}{7.9701}{rm}{49}%
}
\]
由于值域空间的维数为\allowbreak{}\hskip.1em \RB{49.617}{7.770}{1.401}{%
\RG{0.000}{0.004}{10.9091}{mi}{110}\RG{8.973}{0.004}{10.9091}{rm}{43}\RG{19.871}{0.004}{10.9091}{rm}{49}\RG{28.367}{0.004}{10.9091}{mi}{60}\RG{39.866}{0.004}{10.9091}{mi}{109}%
}\hskip.1em\allowbreak{} ，由秩—零度公式，存在不全为零的实数\allowbreak{}\hskip.1em \RB{7.748}{5.326}{2.213}{%
\RG{0.000}{0.004}{10.9091}{mi}{99}\RG{4.721}{-1.633}{7.9701}{mi}{105}%
}\hskip.1em\allowbreak{} ，使

\[
\RB{129.287}{17.656}{13.663}{%
\RG{4.133}{13.641}{7.9701}{mi}{109}\RG{85.148}{13.641}{7.9701}{mi}{109}\RG{0.000}{10.368}{10.9092}{ex}{88}\RG{81.015}{10.368}{10.9092}{ex}{88}\RG{17.576}{0.004}{10.9091}{mi}{99}\RG{25.678}{0.004}{10.9091}{mi}{118}\RG{37.377}{0.004}{10.9091}{rm}{61}\RG{48.876}{0.004}{10.9091}{rm}{48}\RG{54.347}{0.004}{10.9091}{mi}{59}\RG{98.590}{0.004}{10.9091}{mi}{99}\RG{109.723}{0.004}{10.9091}{rm}{61}\RG{121.221}{0.004}{10.9091}{rm}{48}\RG{126.692}{0.004}{10.9091}{mi}{58}\RG{22.297}{-1.632}{7.9701}{mi}{105}\RG{30.966}{-1.632}{7.9701}{mi}{105}\RG{103.311}{-1.632}{7.9701}{mi}{105}\RG{1.027}{-13.083}{7.9701}{mi}{105}\RG{3.910}{-13.083}{7.9701}{rm}{61}\RG{10.493}{-13.083}{7.9701}{rm}{49}\RG{82.041}{-13.083}{7.9701}{mi}{105}\RG{84.925}{-13.083}{7.9701}{rm}{61}\RG{91.508}{-13.083}{7.9701}{rm}{49}%
}
\]
后一等式保证这些非零系数中既有正数，也有负数。记

\[
\RB{182.907}{8.685}{3.760}{%
\RG{0.000}{0.003}{10.9091}{mi}{80}\RG{11.549}{0.003}{10.9091}{rm}{61}\RG{23.064}{0.003}{10.9091}{sy}{102}\RG{28.519}{0.003}{10.9091}{mi}{105}\RG{35.307}{0.003}{10.9091}{rm}{58}\RG{41.368}{0.003}{10.9091}{mi}{99}\RG{52.501}{0.003}{10.9091}{mi}{62}\RG{64.017}{0.003}{10.9091}{rm}{48}\RG{69.471}{0.003}{10.9091}{sy}{103}\RG{74.926}{0.003}{10.9091}{mi}{59}\RG{101.577}{0.003}{10.9091}{mi}{78}\RG{114.577}{0.003}{10.9091}{rm}{61}\RG{126.092}{0.003}{10.9091}{sy}{102}\RG{131.547}{0.003}{10.9091}{mi}{106}\RG{139.694}{0.003}{10.9091}{rm}{58}\RG{145.755}{0.003}{10.9091}{mi}{99}\RG{157.888}{0.003}{10.9091}{mi}{60}\RG{169.403}{0.003}{10.9091}{rm}{48}\RG{174.858}{0.003}{10.9091}{sy}{103}\RG{180.312}{0.003}{10.9091}{mi}{58}\RG{46.089}{-1.634}{7.9701}{mi}{105}\RG{150.476}{-1.634}{7.9701}{mi}{106}%
}
\]
则\allowbreak{}\hskip.1em \RB{22.253}{7.950}{2.606}{%
\RG{0.000}{-0.001}{10.9091}{mi}{80}\RG{7.298}{-0.001}{10.9091}{mi}{59}\RG{12.142}{-0.001}{10.9091}{mi}{78}%
}\hskip.1em\allowbreak{} 均非空，而且

\[
\RB{130.011}{10.867}{15.387}{%
\RG{25.679}{10.367}{10.9092}{ex}{88}\RG{75.244}{10.367}{10.9092}{ex}{88}\RG{0.000}{0.004}{10.9091}{mi}{119}\RG{11.134}{0.004}{10.9091}{rm}{58}\RG{14.156}{0.004}{10.9091}{rm}{61}\RG{43.255}{0.004}{10.9091}{mi}{99}\RG{51.357}{0.004}{10.9091}{mi}{118}\RG{63.058}{0.004}{10.9091}{rm}{61}\RG{91.673}{0.004}{10.9091}{rm}{40}\RG{95.915}{0.004}{10.9091}{sy}{0}\RG{104.400}{0.004}{10.9091}{mi}{99}\RG{113.503}{0.004}{10.9091}{rm}{41}\RG{117.746}{0.004}{10.9091}{mi}{118}\RG{127.416}{0.004}{10.9091}{mi}{58}\RG{47.976}{-1.633}{7.9701}{mi}{105}\RG{56.645}{-1.633}{7.9701}{mi}{105}\RG{109.121}{-1.633}{7.9701}{mi}{106}\RG{123.034}{-1.633}{7.9701}{mi}{106}\RG{26.021}{-13.261}{7.9701}{mi}{105}\RG{28.905}{-13.261}{7.9701}{sy}{50}\RG{34.550}{-13.261}{7.9701}{mi}{80}\RG{74.573}{-13.261}{7.9701}{mi}{106}\RG{78.457}{-13.261}{7.9701}{sy}{50}\RG{84.103}{-13.261}{7.9701}{mi}{78}%
}
\]
分别用这两个表达式计算内积，得到

\[
\RB{163.347}{10.859}{15.395}{%
\RG{38.290}{10.359}{10.9092}{ex}{88}\RG{56.537}{10.359}{10.9092}{ex}{88}\RG{19.012}{4.500}{7.9701}{rm}{50}\RG{0.000}{-0.005}{10.9091}{sy}{107}\RG{5.454}{-0.005}{10.9091}{mi}{119}\RG{13.558}{-0.005}{10.9091}{sy}{107}\RG{26.775}{-0.005}{10.9091}{rm}{61}\RG{74.785}{-0.005}{10.9091}{mi}{99}\RG{82.887}{-0.005}{10.9091}{rm}{40}\RG{87.129}{-0.005}{10.9091}{sy}{0}\RG{95.614}{-0.005}{10.9091}{mi}{99}\RG{104.718}{-0.005}{10.9091}{rm}{41}\RG{108.960}{-0.005}{10.9091}{sy}{104}\RG{113.202}{-0.005}{10.9091}{mi}{118}\RG{121.872}{-0.005}{10.9091}{mi}{59}\RG{126.705}{-0.005}{10.9091}{mi}{118}\RG{136.390}{-0.005}{10.9091}{sy}{105}\RG{143.663}{-0.005}{10.9091}{mi}{60}\RG{155.178}{-0.005}{10.9091}{rm}{48}\RG{160.632}{-0.005}{10.9091}{mi}{59}\RG{79.506}{-1.641}{7.9701}{mi}{105}\RG{100.335}{-1.641}{7.9701}{mi}{106}\RG{118.490}{-1.641}{7.9701}{mi}{105}\RG{132.008}{-1.641}{7.9701}{mi}{106}\RG{38.632}{-13.269}{7.9701}{mi}{105}\RG{41.515}{-13.269}{7.9701}{sy}{50}\RG{47.161}{-13.269}{7.9701}{mi}{80}\RG{55.866}{-13.269}{7.9701}{mi}{106}\RG{59.750}{-13.269}{7.9701}{sy}{50}\RG{65.396}{-13.269}{7.9701}{mi}{78}%
}
\]
因为\allowbreak{}\hskip.1em \RB{53.506}{7.951}{0.740}{%
\RG{0.000}{0.000}{10.9091}{mi}{80}\RG{10.943}{0.000}{10.9091}{sy}{92}\RG{20.640}{0.000}{10.9091}{mi}{78}\RG{33.624}{0.000}{10.9091}{rm}{61}\RG{45.140}{0.000}{10.9091}{bb}{63}%
}\hskip.1em\allowbreak{} ，每项均严格为负，且至少有一项。这与模平方非负矛盾，故\allowbreak{}\hskip.1em \RB{49.065}{7.766}{1.994}{%
\RG{0.000}{0.000}{10.9091}{mi}{109}\RG{12.608}{0.000}{10.9091}{sy}{20}\RG{24.123}{0.000}{10.9091}{mi}{110}\RG{33.095}{0.000}{10.9091}{rm}{43}\RG{43.994}{0.000}{10.9091}{rm}{49}%
}\hskip.1em\allowbreak{} 。

下面给出取等构造。在\allowbreak{}\hskip.1em \RB{23.664}{9.755}{0.504}{%
\RG{0.000}{-0.004}{10.9091}{bb}{82}\RG{7.879}{3.955}{7.9701}{mi}{110}\RG{13.018}{3.955}{7.9701}{rm}{43}\RG{19.601}{3.955}{7.9701}{rm}{49}%
}\hskip.1em\allowbreak{} 中令

\[
\RB{300.734}{19.592}{14.125}{%
\RG{24.500}{18.655}{10.9092}{ex}{40}\RG{97.746}{18.655}{10.9092}{ex}{41}\RG{48.614}{13.637}{7.9701}{mi}{110}\RG{53.752}{13.637}{7.9701}{rm}{43}\RG{60.336}{13.637}{7.9701}{rm}{49}\RG{48.714}{10.364}{10.9092}{ex}{88}\RG{192.122}{7.381}{10.9091}{rm}{49}\RG{0.000}{0.001}{10.9091}{mi}{72}\RG{12.985}{0.001}{10.9091}{rm}{61}\RG{33.288}{0.001}{10.9091}{mi}{120}\RG{42.553}{0.001}{10.9091}{rm}{58}\RG{66.392}{0.001}{10.9091}{mi}{120}\RG{80.777}{0.001}{10.9091}{rm}{61}\RG{92.275}{0.001}{10.9091}{rm}{48}\RG{108.352}{0.001}{10.9091}{mi}{59}\RG{135.003}{0.001}{10.9091}{mi}{117}\RG{147.676}{0.001}{10.9091}{rm}{61}\RG{159.191}{0.001}{10.9091}{mi}{101}\RG{170.076}{0.001}{10.9091}{sy}{0}\RG{208.713}{0.001}{10.9091}{rm}{49}\RG{225.895}{0.001}{10.9091}{rm}{40}\RG{230.128}{0.001}{10.9091}{rm}{49}\RG{238.622}{0.001}{10.9091}{sy}{20}\RG{250.137}{0.001}{10.9091}{mi}{105}\RG{256.926}{0.001}{10.9091}{sy}{20}\RG{268.441}{0.001}{10.9091}{mi}{110}\RG{277.413}{0.001}{10.9091}{rm}{43}\RG{288.312}{0.001}{10.9091}{rm}{49}\RG{293.766}{0.001}{10.9091}{rm}{41}\RG{298.019}{0.001}{10.9091}{mi}{59}\RG{141.264}{-1.636}{7.9701}{mi}{105}\RG{164.270}{-1.636}{7.9701}{mi}{105}\RG{72.627}{-1.777}{7.9701}{mi}{107}\RG{182.181}{-7.483}{10.9091}{mi}{110}\RG{191.154}{-7.482}{10.9091}{rm}{43}\RG{202.052}{-7.482}{10.9091}{rm}{49}\RG{48.872}{-13.352}{7.9701}{mi}{107}\RG{53.494}{-13.352}{7.9701}{rm}{61}\RG{60.078}{-13.352}{7.9701}{rm}{49}\RR{182.181}{2.728}{25.336}{0.436}%
}
\]
其中\allowbreak{}\hskip.1em \RB{63.865}{8.679}{3.230}{%
\RG{0.000}{-0.003}{10.9091}{rm}{49}\RG{9.303}{-0.003}{10.9091}{rm}{61}\RG{20.801}{-0.003}{10.9091}{rm}{40}\RG{25.034}{-0.003}{10.9091}{rm}{49}\RG{30.515}{-0.003}{10.9091}{mi}{59}\RG{35.348}{-0.003}{10.9091}{mi}{58}\RG{40.192}{-0.003}{10.9091}{mi}{58}\RG{45.024}{-0.003}{10.9091}{mi}{58}\RG{49.868}{-0.003}{10.9091}{mi}{59}\RG{54.757}{-0.003}{10.9091}{rm}{49}\RG{60.212}{-0.003}{10.9091}{rm}{41}%
}\hskip.1em\allowbreak{} T。空间\allowbreak{}\hskip.1em \RB{10.111}{7.948}{0.503}{%
\RG{0.000}{-0.003}{10.9091}{mi}{72}%
}\hskip.1em\allowbreak{} 的维数为\allowbreak{}\hskip.1em \RB{6.729}{5.319}{0.623}{%
\RG{0.000}{-0.003}{10.9091}{mi}{110}%
}\hskip.1em\allowbreak{} ，并且\allowbreak{}\hskip.1em \RB{33.070}{7.948}{2.219}{%
\RG{0.000}{-0.003}{10.9091}{mi}{117}\RG{6.245}{-1.639}{7.9701}{mi}{105}\RG{12.656}{-0.003}{10.9091}{sy}{50}\RG{22.959}{-0.003}{10.9091}{mi}{72}%
}\hskip.1em\allowbreak{} 。直接计算可得

\[
\RB{263.962}{15.150}{8.885}{%
\RG{50.403}{7.384}{10.9091}{mi}{110}\RG{181.961}{7.384}{10.9091}{rm}{49}\RG{20.536}{4.509}{7.9701}{rm}{50}\RG{0.000}{0.004}{10.9091}{sy}{107}\RG{5.455}{0.004}{10.9091}{mi}{117}\RG{15.081}{0.004}{10.9091}{sy}{107}\RG{28.298}{0.004}{10.9091}{rm}{61}\RG{70.571}{0.004}{10.9091}{mi}{62}\RG{82.087}{0.004}{10.9091}{rm}{48}\RG{87.541}{0.004}{10.9091}{mi}{59}\RG{114.208}{0.004}{10.9091}{sy}{104}\RG{118.450}{0.004}{10.9091}{mi}{117}\RG{128.076}{0.004}{10.9091}{mi}{59}\RG{132.909}{0.004}{10.9091}{mi}{117}\RG{143.552}{0.004}{10.9091}{sy}{105}\RG{150.825}{0.004}{10.9091}{rm}{61}\RG{162.340}{0.004}{10.9091}{sy}{0}\RG{201.583}{0.004}{10.9091}{mi}{60}\RG{213.098}{0.004}{10.9091}{rm}{48}\RG{229.462}{0.004}{10.9091}{rm}{40}\RG{233.704}{0.004}{10.9091}{mi}{105}\RG{240.493}{0.004}{10.9091}{sy}{54}\RG{240.493}{0.004}{10.9091}{rm}{61}\RG{252.008}{0.004}{10.9091}{mi}{106}\RG{257.125}{0.004}{10.9091}{rm}{41}\RG{261.367}{0.004}{10.9091}{mi}{58}\RG{11.700}{-1.632}{7.9701}{mi}{105}\RG{124.695}{-1.632}{7.9701}{mi}{105}\RG{139.170}{-1.632}{7.9701}{mi}{106}\RG{41.009}{-7.479}{10.9091}{mi}{110}\RG{49.981}{-7.479}{10.9091}{rm}{43}\RG{60.880}{-7.479}{10.9091}{rm}{49}\RG{172.020}{-7.479}{10.9091}{mi}{110}\RG{180.993}{-7.479}{10.9091}{rm}{43}\RG{191.892}{-7.479}{10.9091}{rm}{49}\RR{41.009}{2.732}{25.336}{0.436}\RR{172.020}{2.732}{25.336}{0.436}%
}
\]
因此这\allowbreak{}\hskip.1em \RB{24.642}{7.762}{1.408}{%
\RG{0.000}{-0.003}{10.9091}{mi}{110}\RG{8.825}{-0.003}{10.9091}{rm}{43}\RG{19.571}{-0.003}{10.9091}{rm}{49}%
}\hskip.1em\allowbreak{} 个向量两两夹角大于直角，且夹角的余弦均为\allowbreak{}\hskip.1em \RB{26.123}{8.679}{3.230}{%
\RG{0.000}{-0.003}{10.9091}{sy}{0}\RG{8.485}{-0.003}{10.9091}{rm}{49}\RG{13.940}{-0.003}{10.9091}{rm}{47}\RG{19.394}{-0.003}{10.9091}{mi}{110}%
}\hskip.1em\allowbreak{} 。在\allowbreak{}\hskip.1em \RB{10.111}{7.948}{0.503}{%
\RG{0.000}{-0.003}{10.9091}{mi}{72}%
}\hskip.1em\allowbreak{} 和\allowbreak{}\hskip.1em \RB{8.889}{7.948}{0.743}{%
\RG{0.000}{-0.003}{10.9091}{mi}{86}%
}\hskip.1em\allowbreak{} 中各取一组标准正交基，将前者映到后者便得到保内积的线性同构，故同样的构造适用于\allowbreak{}\hskip.1em \RB{8.889}{7.953}{0.738}{%
\RG{0.000}{0.002}{10.9091}{mi}{86}%
}\hskip.1em\allowbreak{} 。\allowbreak{}\hskip.1em \RB{8.366}{7.997}{0.500}{%
\RG{0.000}{0.002}{10.9091}{am}{3}%
}\hskip.1em\allowbreak{}

方法小结：增加系数和为零的条件，使线性关系必含正、负两类系数，再以模平方导出矛盾。取等构造来自正单纯形各顶点减去其重心。

题目来源：谢启鸿《复旦大学高等代数习题集（2023版）》，第117页。


\end{document}
